\documentclass[11pt]{amsart}

\usepackage[T1]{fontenc}
\usepackage[utf8]{inputenc}
\usepackage{lmodern}
\usepackage{microtype}
\usepackage{mathtools,amssymb,amsfonts,amsthm}
\usepackage{mathrsfs}
\usepackage{geometry}
\usepackage{booktabs}
\usepackage{listings}
\usepackage{xcolor}
\lstset{
 basicstyle=\ttfamily\footnotesize,
 breaklines=true,
 columns=fullflexible,
 frame=single,
 backgroundcolor=\color{gray!5}
}
\usepackage[hidelinks]{hyperref}

\geometry{margin=1in}
\numberwithin{equation}{section}
\setcounter{tocdepth}{3}
\setcounter{secnumdepth}{3}

\newtheorem{theorem}{Theorem}[section]
\newtheorem{proposition}[theorem]{Proposition}
\newtheorem{corollary}[theorem]{Corollary}
\newtheorem{lemma}[theorem]{Lemma}
\newtheorem{definition}[theorem]{Definition}
\newtheorem{remark}[theorem]{Remark}
\newtheorem{example}[theorem]{Example}

\newcommand{\R}{\mathbb R}
\newcommand{\C}{\mathbb C}
\newcommand{\Hcal}{\mathcal H}
\newcommand{\Pcal}{\mathcal P}
\newcommand{\dd}{\,d}
\newcommand{\ip}[2]{\left\langle #1,#2\right\rangle}
\newcommand{\norm}[1]{\left\lVert #1\right\rVert}
\newcommand{\PP}{\mathbb P}
\newcommand{\Dq}{D_q}

\title[Critical-Parameter Rigidity for the Riemann Zeta Function]{Critical-Parameter Rigidity for the Riemann Zeta Function under $L^2$-Projective Coincidence}
\author{Rotkiv Aynäsac}
\address{}
\thanks{}
\date{Version 14 --- August 2026}

\begin{document}

\begin{abstract}
For a fixed interval between consecutive critical-line zeros of the Riemann
zeta function, we use the Riemann--Siegel phase to transport every vertical
fibre $\zeta(\beta+it)$, $1/2\leq\beta\leq1$, to one common parameter
interval.  In the Hilbert space $L^2([0,1])$ we remove the four-dimensional
polynomial subspace of degree at most three and then projectivize.  The main
unconditional result is a rigidity theorem: two projected vertical zeta
fibres have the same projective class if and only if their real parameters
are equal.  In particular, coincidence with the class at $\beta=1/2$ forces
$\beta=1/2$.  The proof is global and elementary once the phase coordinate
has been fixed: a hypothetical projective coincidence produces a
finite-dimensional functional identity, analytic continuation transports it
to the simple pole of $\zeta$ at $s=1$, and the pole cannot be cancelled by
the remaining terms.

We then separate this genuine zeta-specific rigidity from the universal
holomorphic transport equation.  The entire function $g(s)=s-s_0$ satisfies
exactly the same Cauchy--Riemann transport law, conformality, positive
Jacobian and harmonicity, while possessing a zero at any prescribed
$\beta_0\neq1/2$.  Thus these universal geometric identities cannot by
themselves confine zeros to the critical line.  The missing zero-specific
step is isolated as an exact projective-collapse condition.  We prove that
this condition is equivalent to the Riemann hypothesis, and also give
finite-dimensional, fourth-derivative and Wronskian formulations of the same
condition.  Finally, a quantitative version based on the Fubini--Study defect
shows how a zero-triggered \emph{almost} coincidence, if established by a
zeta-specific arithmetic estimate, would force quantitative proximity to the
critical line without assuming exact collapse.
\end{abstract}

\maketitle
\tableofcontents
\bigskip

\section{Introduction}
\label{sec:introduction}

Write
\begin{equation}
 s=\beta+it,
 \qquad
 \frac12\leq\beta\leq1,
 \label{eq:s-beta-t}
\end{equation}
and consider the vertical family, in the setting initiated by Riemann's original memoir \cite{Riemann1859},
\begin{equation}
 \boxed{\zeta(\beta+it).}
 \label{eq:vertical-zeta-family}
\end{equation}
The object studied in this paper is the Riemann zeta function itself.  No
auxiliary replacement zeta function is introduced.  On each interval between
consecutive critical-line zeros, the Riemann--Siegel phase supplies a common
coordinate; the corresponding vertical restrictions are then placed in the
Hilbert space $L^2([0,1])$, quotiented by the four-dimensional polynomial
space of degree at most three, and finally projectivized.  All three operations
are fixed before the real parameter $\beta$ is varied.

The central unconditional result is the global rigidity statement
\begin{equation}
\boxed{
 [Q_\perp\zeta(\alpha+i t_n(\cdot))]
 =
 [Q_\perp\zeta(\beta+i t_n(\cdot))]
 \quad\Longleftrightarrow\quad
 \alpha=\beta,
}
\label{eq:intro-global-rigidity}
\end{equation}
valid on every fixed phase cell.  In particular,
\begin{equation}
\boxed{
 [Q_\perp\zeta(\beta+i t_n(\cdot))]
 =
 [Q_\perp\zeta(\tfrac12+i t_n(\cdot))]
 \quad\Longleftrightarrow\quad
 \beta=\frac12.
}
\label{eq:intro-critical-rigidity}
\end{equation}
The proof is exact rather than asymptotic.  Projective coincidence produces a
finite-dimensional identity; analytic continuation carries that identity to
the simple pole of $\zeta$ at $s=1$; and the pole cannot be cancelled by the
remaining fibre or by the discarded cubic component.

\subsection{Where this construction sits among classical RH criteria}
\label{subsec:intro-literature-position}

The present construction lies at the intersection of two old and substantial
traditions, but is not a restatement of either one.  The first is the
Hilbert-space tradition surrounding the Nyman--Beurling criterion and its
arithmetic refinements: RH can be encoded as an $L^2$ closure problem for
fractional-part dilates \cite{Nyman1950,Beurling1955,BaezDuarte2000,BaezDuarte2002,
BaezDuarte2003,Burnol2002}.  The second is the complex-analytic rigidity tradition going
back to Ostrowski and Voronin and continuing through modern differential,
functional, and difference-independence results for $\zeta$
\cite{Ostrowski1920,Voronin1973,KaczorowskiPerelli1998,ChiangFeng2006,Lu2024,
KorolevLaurincikas2025}.  Here the Hilbert-space quotient is local in a
Riemann--Siegel phase cell, whereas the decisive uniqueness mechanism is
meromorphic: the pole at $s=1$ separates every two distinct vertical
parameters after the quotient.

A closely related older line is the geometric/derivative route represented by
Speiser's criterion and the quantitative work of Levinson--Montgomery on
zeros of derivatives of $\zeta$ \cite{Speiser1934,LevinsonMontgomery1974}.
Those results translate information about the zero set into restrictions on
$\zeta'$ and its derivatives.  The present construction does something
different but complementary: differentiation is taken only after a fixed
Riemann--Siegel phase reparametrization, the polynomial components of
degree at most three are removed in $L^2$, and uniqueness is then tested
projectively.  Thus the paper
combines two classical ideas---Hilbert-space reduction and analytic
zero/derivative rigidity---without reducing to either classical criterion.

This should also be distinguished from the global positivity/equivalent-
criterion tradition of Weil and Li \cite{Weil1952,Li1997,BombieriLagarias1999}.
Those criteria encode all zeros through global test functions, explicit
formulae, or positivity sequences.  The present criterion instead compares
two vertical zeta fibres on one fixed cell and asks whether their classes in a
finite-codimensional projective Hilbert quotient can coincide.  Within the
literature reviewed here, we have not found a criterion with the simultaneous
combination
\begin{equation}
 \boxed{
 \text{Riemann--Siegel cell}
 +L^2/\Pcal_3
 +\text{projective coincidence}
 +\text{pole separation at }s=1.
 }
 \label{eq:intro-architecture}
\end{equation}
The components are classical; the way they are coupled in
\eqref{eq:intro-architecture}, and the resulting exact critical-parameter
rigidity, are the distinctive content of the construction.

The unconditional theorem has value independently of the unresolved RH
trigger.  It identifies a concrete one-parameter family in an infinite-
dimensional Hilbert space for which quotienting by a fixed four-dimensional
subspace and then forgetting scalar magnitude still loses no information
about the horizontal parameter.  That injectivity is not a generic Hilbert-
space fact: it is enforced here by the global meromorphic singularity of
$\zeta$.  In particular, the construction yields an exact bridge from local
$L^2$ projective geometry to global analytic continuation, rather than an
asymptotic numerical separation.  The RH-equivalent condition introduced
later should therefore be read as an additional arithmetic trigger attached
to an already nontrivial rigidity theorem, not as the source of the theorem's
content.

There is a second neighboring literature that must be kept logically
separate.  Values of $\zeta$ or its derivatives at shifted zero ordinates have
been studied through discrete mean values, universality, and local value
distribution
\cite{Gonek1984,GarunkstisLaurincikas2018,
GarunkstisLaurincikasMacaitiene2017,Laurincikas2019,Cicek2021,
LaurincikasSiauciunas2020,GarunkstisNovikas2025,GarunkstisPaliulionyte2025}.  These results contain rich
information about $\zeta(\beta+i\gamma)$ and nearby shifts, but none of the
cited theorems turns the scalar event $\zeta(\beta+i\gamma)=0$ into the exact
whole-cell finite-dimensional relation used below.  Difference-independence,
on the other hand, supports the rigidity side of the paper, not the missing
zero-trigger.

Finite-dimensional translation-invariant spaces supply another relevant
background: such spaces are governed by exponential-polynomial structure
\cite{AnseloneKorevaar1964,Engert1970,Szekelyhidi1982}.  This makes a
finite-rank collapse a natural object to test, but it does not make a point
zero force that collapse.  The distinction between \emph{classification of a
finite-dimensional relation} and \emph{generation of that relation by a
zero} is one of the main logical separations of this paper.

\subsection{Why the route is not a zero-free-region approach to \texorpdfstring{$1/2$}{1/2}}
\label{subsec:intro-no-log-barrier}

A classical route to RH proceeds through prime-number-theorem error terms or
zero-free regions.  Von Koch connected RH with square-root-scale prime
counting errors \cite{vonKoch1901}; modern unconditional zero-free regions
remain of de la Vall\'ee Poussin or Vinogradov--Korobov type and approach
$\Re s=1$ rather than $\Re s=1/2$
\cite{Ford2002,MossinghoffTrudgianYang2024}.  In that framework, improving an
error exponent all the way to the critical exponent $1/2$ is not an
innocuous limiting process: at the relevant uniform strength it is itself an
RH-level statement.  A hypothetical zero with real part $\beta>1/2$ produces
an $x^\beta$ contribution in the explicit formula and cannot remain uniformly
hidden below a square-root-scale error.

The present route does not attempt to move a contour through a sequence of
zero-free strips, nor does it require a bound of the form
$\beta\leq1/2+\delta(T)$ with $\delta(T)$ shrinking logarithmically.  Its
unconditional separation is exact for every fixed $\beta>1/2$:
\begin{equation}
 M_n(\beta)>0
 \qquad(\beta>1/2),
 \label{eq:intro-exact-separation-all-beta}
\end{equation}
while at the same time
\begin{equation}
 M_n(\beta)
 =P_n(1/2)(\beta-1/2)^2
 +o\!\left((\beta-1/2)^2\right)
 \qquad(\beta\downarrow1/2).
 \label{eq:intro-quadratic-approach}
\end{equation}
Thus the projective classes may approach the critical class arbitrarily
closely, and in fact do so quadratically, but they cannot coincide at any
$\beta>1/2$.  No uniform positive gap from $1/2$ is claimed or needed.  This
is precisely why the usual ``how close can one push a zero-free boundary to
$1/2$?'' obstruction does not condemn the rigidity theorem: the theorem is a
nonperturbative uniqueness statement, not a zero-free-region estimate.

This does not remove the difficulty of RH; it relocates it.  The projective
side is closed, but the missing implication is
\begin{equation}
 \zeta(\beta+i\gamma)=0
 \quad\Longrightarrow\quad
 [Q_\perp\zeta(\beta+i t_n(\cdot))]
 =
 [Q_\perp\zeta(\tfrac12+i t_n(\cdot))].
 \label{eq:intro-zero-collapse}
\end{equation}
Theorem~\ref{thm:RH-equivalence} proves that
\eqref{eq:intro-zero-collapse} is equivalent to RH once the unconditional
rigidity is available.  It is therefore not presented as a small technical
lemma.  What the construction achieves is a sharp localization of the RH
content: everything to the right of the zero-trigger is rigid and explicit;
what remains is to explain why the arithmetic event ``$\zeta$ has a zero''
should force exact or quantitative projective collapse.

\subsection{A possible intermediate target: almost-coincidence}
\label{subsec:intro-almost-collapse}

The exact implication \eqref{eq:intro-zero-collapse} is RH-equivalent, but the
geometry also produces a weaker target that is not tautologically RH.  If one
can prove, at a zero,
\begin{equation}
 M_n(\beta)\leq\varepsilon_n(\gamma),
 \label{eq:intro-upper-defect}
\end{equation}
then the already proved lower rigidity converts this into a bound for
$|\beta-1/2|$.  Locally,
\begin{equation}
 \left|\beta-\frac12\right|
 \leq
 \sqrt{\frac{2\varepsilon_n(\gamma)}{P_n(1/2)}}
 \label{eq:intro-proximity}
\end{equation}
under the explicit threshold stated in
Theorem~\ref{thm:almost-collapse-proximity}.  This gives the construction
independent value even before exact collapse: a zeta-specific arithmetic
upper bound for the projective defect would immediately become a quantitative
horizontal bound for zeros.  It is a different object from a smooth fit to
zero ordinates and from a classical zero-free strip.

Finally, the universal Cauchy--Riemann transport must not be confused with the
zeta-specific rigidity.  Section~\ref{sec:universal-obstruction} gives the
entire counterexample $g(s)=s-s_0$: it obeys the same transport equation,
conformality, positive Jacobian, and harmonicity, while having a zero at any
prescribed $\beta_0\neq1/2$.  Hence no proof of
\eqref{eq:intro-zero-collapse} can be built from those universal geometric
identities alone.  The role of the pole at $s=1$ is therefore twofold: it
proves the exact projective uniqueness, and it marks the first genuinely
zeta-specific ingredient in the route.

\section{Analytic setup and the common phase coordinate}
\label{sec:phase-setup}

We use standard facts about $\zeta$, Hardy's real $Z$-function, the
Riemann--Siegel phase and the functional equation as in
\cite{Titchmarsh1986,Ivic1985,Edwards1974}; special-function identities for
$\Gamma$, $\psi$ and their derivatives are also recorded in the NIST DLMF
\cite{DLMF}.  Define
\begin{equation}
 \vartheta(t)
 =\Im\log\Gamma\!\left(\frac14+\frac{it}{2}\right)
 -\frac t2\log\pi,
 \label{eq:theta-definition}
\end{equation}
with the standard continuous branch on the positive real axis.  Then
\begin{equation}
 \zeta\!\left(\frac12+it\right)
 =Z(t)e^{-i\vartheta(t)},
 \qquad
 Z(t)\in\R.
 \label{eq:hardy-representation}
\end{equation}

Let
\begin{equation}
 0<\gamma_1<\gamma_2<\cdots
 \label{eq:critical-zero-ordinates}
\end{equation}
be the increasing sequence of distinct positive ordinates for which
$\zeta(1/2+i\gamma_n)=0$.  Hardy's theorem gives infinitely many such
ordinates \cite{Hardy1914}; discreteness of zeros implies that they have no finite
accumulation point \cite{Titchmarsh1986}.

\begin{lemma}[Positivity of the phase derivative on every positive zero cell]
\label{lem:theta-positive}
For every $n\geq1$,
\begin{equation}
 \vartheta'(t)>0
 \qquad
 (\gamma_n\leq t\leq\gamma_{n+1}).
 \label{eq:theta-positive}
\end{equation}
Consequently
\begin{equation}
 \Omega_n:=\vartheta(\gamma_{n+1})-\vartheta(\gamma_n)>0.
 \label{eq:omega-positive}
\end{equation}
\end{lemma}

\begin{proof}
Differentiating \eqref{eq:theta-definition} gives
\begin{align}
 \vartheta'(t)
 &=\frac12\Re\psi\!\left(\frac14+\frac{it}{2}\right)
   -\frac12\log\pi,
 \label{eq:theta-digamma}\\
 \vartheta''(t)
 &=-\frac14\Im\psi'\!\left(\frac14+\frac{it}{2}\right).
 \label{eq:theta-trigamma}
\end{align}
For $\Re z>0$ the trigamma series
\begin{equation}
 \psi'(z)=\sum_{k=0}^{\infty}\frac1{(k+z)^2}
 \label{eq:trigamma-series}
\end{equation}
is absolutely convergent.  With $z=1/4+it/2$ and $t>0$,
\begin{equation}
 \Im\frac1{(k+\frac14+\frac{it}{2})^2}
 =-
 \frac{(k+\frac14)t}
 {((k+\frac14)^2+t^2/4)^2}<0.
 \label{eq:trigamma-imag-negative}
\end{equation}
Absolute convergence permits termwise summation, so
$\vartheta''(t)>0$ for $t>0$.  Hence $\vartheta'$ is strictly increasing.

A rigorous bound follows by truncating
\begin{equation}
 \psi(z)=-\gamma_E+\sum_{k=0}^{\infty}
 \left(\frac1{k+1}-\frac1{k+z}\right)
 \label{eq:digamma-series}
\end{equation}
at $k=199$. With $z=\tfrac14+7i$ (i.e.\ $t=14$), write
\begin{equation}
 \Re\frac1{k+z}=\frac{k+\tfrac14}{(k+\tfrac14)^2+49},
\end{equation}
so that the partial sum
\begin{equation}
 S_{199}:=\sum_{k=0}^{199}
 \left(\frac1{k+1}-\frac{k+\tfrac14}{(k+\tfrac14)^2+49}\right)
 =2.526047\ldots
 \label{eq:S199-value}
\end{equation}
is an explicit finite rational-arithmetic quantity. For the tail,
$\bigl|\Re\bigl[\tfrac1{k+1}-\tfrac1{k+z}\bigr]\bigr|\leq\tfrac34/(k+\tfrac14)$
for $k\geq200$ (both terms are $O(1/k)$ with the stated constant once
$k\geq200>49$), giving
\begin{equation}
 \left|\sum_{k=200}^\infty
 \Re\left(\frac1{k+1}-\frac1{k+z}\right)\right|
 \leq\frac34\sum_{k=200}^\infty\frac1{(k+\tfrac14)^2}
 \leq\frac3{4\cdot199}=0.0037688\ldots.
 \label{eq:tail-bound-explicit}
\end{equation}
Combining \eqref{eq:S199-value} and \eqref{eq:tail-bound-explicit} with
$\gamma_E=0.5772156\ldots$ gives the certified interval
\begin{equation}
 \Re\psi\!\left(\tfrac14+7i\right)
 \in[1.945063,\,1.952600],
 \label{eq:Repsi-certified}
\end{equation}
and hence, from \eqref{eq:theta-digamma} with $\log\pi=1.144730\ldots$,
\begin{equation}
 \vartheta'(14)\in[0.400166,\,0.403935]\subset(0.40016,\infty).
 \label{eq:theta-prime-14-certified}
\end{equation}
This certified interval contains the independently computed reference value
$\vartheta'(14)=0.4004837\ldots$, confirming the bound.  The first positive critical-line zero satisfies
$\gamma_1=14.134725\ldots>14$; this is also contained in rigorous
computational verifications of the critical line, far beyond this height
\cite{PlattTrudgian2021}.  Since $\vartheta'$ is increasing,
\eqref{eq:theta-positive} follows for every $t\geq\gamma_1$, and then
\eqref{eq:omega-positive} follows by integration.
\end{proof}

\begin{definition}[Common phase coordinate]
\label{def:phase-coordinate}
For a fixed $n$, define $t_n:[0,1]\to[\gamma_n,\gamma_{n+1}]$ by
\begin{equation}
 \boxed{
 \vartheta(t_n(q))
 =\vartheta(\gamma_n)+\Omega_n q.
 }
 \label{eq:tn-definition}
\end{equation}
Then $t_n$ is a real-analytic increasing diffeomorphism and
\begin{equation}
 \boxed{
 t_n'(q)=\frac{\Omega_n}{\vartheta'(t_n(q))}>0.
 }
 \label{eq:tn-derivative}
\end{equation}
For $1/2\leq\beta\leq1$ we write, only as a shorthand for the same zeta
function,
\begin{equation}
 \zeta_{\beta,n}(q)
 :=\zeta(\beta+i t_n(q)).
 \label{eq:zeta-beta-n}
\end{equation}
\end{definition}

\begin{proof}
Lemma~\ref{lem:theta-positive} and the inverse function theorem give a unique
real-analytic inverse of $\vartheta$ on the compact interval.  Differentiating
\eqref{eq:tn-definition} gives \eqref{eq:tn-derivative}.  Since
$t_n([0,1])\subset(0,\infty)$, none of the points
$\beta+i t_n(q)$ equals the pole $s=1$, so
\eqref{eq:zeta-beta-n} is real-analytic in $q$ and smooth jointly in
$(\beta,q)$ on $[1/2,1]\times[0,1]$.
\end{proof}

\section{Universal holomorphic transport and its geometric consequences}
\label{sec:universal-transport}

The following identities are universal.  They do not use the Euler product,
functional equation, zero distribution, or pole structure of $\zeta$.

\begin{theorem}[Universal quarter-turn geometry]
\label{thm:universal-quarter-turn}
Let $g$ be holomorphic on a domain containing a rectangle in the variable
$s=\beta+it$, and define $G(\beta,t)=g(\beta+it)$.  Then
\begin{equation}
 \boxed{G_\beta=-iG_t.}
 \label{eq:universal-cr}
\end{equation}
At every point,
\begin{align}
 |G_\beta|&=|G_t|,
 \label{eq:universal-isometry}\\
 \Re(\overline{G_\beta}G_t)&=0,
 \label{eq:universal-orthogonality}\\
 J_G&=|G_t|^2=|g'(\beta+it)|^2\geq0,
 \label{eq:universal-jacobian}\\
 G_{\beta\beta}+G_{tt}&=0.
 \label{eq:universal-harmonicity}
\end{align}
Where $g'\neq0$, the pullback of the Euclidean metric is
\begin{equation}
 \boxed{
 g^*(|dw|^2)=|g'(\beta+it)|^2(\dd\beta^2+\dd t^2).
 }
 \label{eq:universal-conformal-metric}
\end{equation}
\end{theorem}

\begin{proof}
Holomorphy gives
$G_\beta=g'(\beta+it)$ and $G_t=i g'(\beta+it)$, proving
\eqref{eq:universal-cr}.  The norm and orthogonality identities follow
immediately.  The real Jacobian of a holomorphic map is $|g'|^2$, and
\eqref{eq:universal-harmonicity} follows by differentiating
\eqref{eq:universal-cr}.  The metric identity follows by evaluating the
squared norm of
$g'(\beta+it)(\dd\beta+i\dd t)$.
\end{proof}

\begin{corollary}[Exact transport in the fixed phase coordinate]
\label{cor:phase-transport}
For any such $g$, put
\begin{equation}
 G_n(\beta,q):=g(\beta+i t_n(q)).
 \label{eq:Gn-definition}
\end{equation}
Then
\begin{equation}
 \boxed{
 \partial_\beta G_n(\beta,q)
 =-i a_n(q)\partial_qG_n(\beta,q),
 \qquad
 a_n(q):=\frac1{t_n'(q)}
 =\frac{\vartheta'(t_n(q))}{\Omega_n}.
 }
 \label{eq:phase-transport}
\end{equation}
\end{corollary}

\begin{proof}
The chain rule gives
$\partial_qG_n=i t_n'(q)g'(\beta+i t_n(q))$ and
$\partial_\beta G_n=g'(\beta+i t_n(q))$.  Since $t_n'(q)>0$,
division by $t_n'(q)$ is legitimate and yields
\eqref{eq:phase-transport}.
\end{proof}

\section{The L2 quotient and projective defect}
\label{sec:l2-quotient}

Let
\begin{equation}
 \Hcal=L^2([0,1],\dd q;\C),
 \qquad
 \ip{f}{h}=\int_0^1\overline{f(q)}h(q)\dd q,
 \label{eq:Hilbert-space}
\end{equation}
with inner product linear in the second variable.  Let
\begin{equation}
 \Pcal_3=\operatorname{span}\{1,q,q^2,q^3\},
 \qquad
 Q_\perp=I-\Pi_3,
 \label{eq:Qperp-definition}
\end{equation}
where $\Pi_3$ is the orthogonal projection onto $\Pcal_3$.

\begin{lemma}[Quotient interpretation]
\label{lem:quotient-interpretation}
The range of $Q_\perp$ is $\Pcal_3^\perp$ and its kernel is $\Pcal_3$.
For $f,h\in\Hcal$,
\begin{equation}
 \ip{[f]}{[h]}_{\Hcal/\Pcal_3}
 :=\ip{Q_\perp f}{Q_\perp h}
 \label{eq:quotient-inner-product}
\end{equation}
defines the canonical Hilbert-space structure on $\Hcal/\Pcal_3$.
\end{lemma}

\begin{proof}
This is the orthogonal decomposition
$\Hcal=\Pcal_3\oplus\Pcal_3^\perp$.  Each quotient class has exactly one
representative in $\Pcal_3^\perp$, namely $Q_\perp f$, and
\eqref{eq:quotient-inner-product} is independent of the chosen representative.
\end{proof}

For a nonzero $u\in\Hcal$, let $[u]\in\PP(\Hcal)$ denote its complex
projective class.  Once nonvanishing has been proved, define
\begin{equation}
 u_{\beta,n}:=Q_\perp\zeta_{\beta,n}
 \label{eq:u-beta-n}
\end{equation}
and the Fubini--Study defect from the critical fibre by
\begin{equation}
 \boxed{
 M_n(\beta)
 =1-
 \frac{|\ip{u_{1/2,n}}{u_{\beta,n}}|^2}
 {\norm{u_{1/2,n}}^2\norm{u_{\beta,n}}^2}.
 }
 \label{eq:M-definition}
\end{equation}
Thus $0\leq M_n(\beta)\leq1$, and $M_n(\beta)=0$ exactly when the two
nonzero vectors have the same projective class.

\section{Holomorphic continuation of the phase coordinate}
\label{sec:phase-holomorphic}

The pole-separation argument requires a complex neighborhood of the real phase
cell.  Define
\begin{equation}
 \Theta(z)=\frac{1}{2i}\left[
 \log\Gamma\!\left(\frac14+\frac{iz}{2}\right)
 -\log\Gamma\!\left(\frac14-\frac{iz}{2}\right)
 \right]-\frac z2\log\pi
 \label{eq:Theta-complex}
\end{equation}
on
\begin{equation}
 S:=\{z\in\C:|\Im z|<1/2\}.
 \label{eq:S-strip}
\end{equation}
Both Gamma arguments have positive real part on $S$.  Since $\Gamma$ has no
zeros and no poles in the open right half-plane, holomorphic logarithms may be
chosen there; these choices can be normalized so that
$\Theta(t)=\vartheta(t)$ for real $t$.

Define
\begin{equation}
 q_n(z)=\frac{\Theta(z)-\Theta(\gamma_n)}{\Omega_n}.
 \label{eq:q-n-holomorphic}
\end{equation}
Then $q_n$ is holomorphic on $S$ and equals the real phase coordinate on the
cell.

\begin{lemma}[Boundary growth at the critical pole]
\label{lem:q-boundary-growth}
Let $z_0=-i/2$.  As $z\to z_0$ from within $S$,
\begin{equation}
 q_n(z)=O\!\left(1+|\log(z-z_0)|\right).
 \label{eq:q-boundary-growth}
\end{equation}
Consequently, for every $p_3\in\Pcal_3$,
\begin{equation}
 p_3(q_n(z))
 =O\!\left(1+|\log(z-z_0)|^3\right).
 \label{eq:p3-boundary-growth}
\end{equation}
\end{lemma}

\begin{proof}
At $z=z_0$, the factor
$1/4-iz/2$ tends to $0$, whereas $1/4+iz/2$ tends to $1/2$.
In the right half-plane, the Laurent expansion
$\Gamma(w)=w^{-1}+O(1)$ implies
$\log\Gamma(w)=-\log w+O(w)$ on a fixed local branch.  Therefore
\eqref{eq:Theta-complex} has at most logarithmic growth at $z_0$, which gives
\eqref{eq:q-boundary-growth}.  A polynomial of degree at most three then
satisfies \eqref{eq:p3-boundary-growth}.
\end{proof}

\section{Pole separation and L2-projective rigidity}
\label{sec:pole-rigidity}

\begin{lemma}[Pole separation]
\label{lem:pole-separation}
Let
\begin{equation}
 \frac12\leq\alpha<\beta\leq1,
 \qquad
 \lambda\in\C^\times,
 \qquad
 p_3\in\Pcal_3.
 \label{eq:pole-separation-hypotheses}
\end{equation}
Then the identity
\begin{equation}
 \zeta(\alpha+i t_n(q))
 -\lambda\zeta(\beta+i t_n(q))
 =p_3(q)
 \label{eq:pole-separation-identity}
\end{equation}
cannot hold on any nonempty real subinterval of $[0,1]$.
\end{lemma}

\begin{proof}
If \eqref{eq:pole-separation-identity} holds in $L^2$ on an interval, the
difference of the two sides is continuous and vanishes almost everywhere.
If it were nonzero at one point, continuity would make it nonzero on an open
subinterval of positive measure, a contradiction.  Thus the identity is
pointwise on the interval.

Replacing $q$ by $q_n(z)$ on the corresponding real $z$-interval gives
\begin{equation}
 \zeta(\alpha+iz)
 -\lambda\zeta(\beta+iz)
 =p_3(q_n(z)).
 \label{eq:pole-separation-complex-identity}
\end{equation}
Both sides are meromorphic on the connected strip $S$.  The identity theorem
for meromorphic functions \cite{Ahlfors1979} therefore extends
\eqref{eq:pole-separation-complex-identity} throughout $S$.

Set
\begin{equation}
 z_\beta=-i(1-\beta).
 \label{eq:z-beta-pole}
\end{equation}
Because $\beta>\alpha\geq1/2$, one has $\beta>1/2$ and hence
$z_\beta\in S$.  Moreover
\begin{equation}
 \beta+iz_\beta=1,
 \qquad
 \alpha+iz_\beta=1-(\beta-\alpha)\neq1.
 \label{eq:pole-values}
\end{equation}
Thus the second zeta term has a simple pole at $z_\beta$, while the first zeta
term and $p_3(q_n(z))$ are holomorphic there.  Since $\lambda\neq0$, the pole
cannot cancel.  This contradiction proves the lemma.
\end{proof}

\begin{proposition}[Nonvanishing of every projected zeta fibre]
\label{prop:projected-nonzero}
For every $n\geq1$ and every $1/2\leq\beta\leq1$,
\begin{equation}
 \boxed{
 Q_\perp\zeta(\beta+i t_n(\cdot))\neq0
 \quad\text{in }L^2([0,1]).
 }
 \label{eq:projected-nonzero}
\end{equation}
\end{proposition}

\begin{proof}
Suppose first that $\beta>1/2$ and
$Q_\perp\zeta_{\beta,n}=0$.  Then
$\zeta_{\beta,n}=p_3$ in $L^2$ for some $p_3\in\Pcal_3$.
Continuity makes the identity pointwise.  Analytic continuation gives
\begin{equation}
 \zeta(\beta+iz)=p_3(q_n(z))
 \label{eq:nonzero-complex-beta}
\end{equation}
on $S$.  At $z=-i(1-\beta)\in S$ the left side has the simple pole of
$\zeta$ at $1$, while the right side is holomorphic, a contradiction.

Now let $\beta=1/2$.  The same assumption and identity theorem give
\begin{equation}
 \zeta\!\left(\frac12+iz\right)=p_3(q_n(z))
 \label{eq:nonzero-complex-critical}
\end{equation}
on $S$.  Write $z=z_0+w$ with $z_0=-i/2$.  Then
\begin{equation}
 \zeta\!\left(\frac12+i(z_0+w)\right)
 =\zeta(1+iw)=\frac1{iw}+O(1),
 \label{eq:zeta-critical-boundary-pole}
\end{equation}
whereas Lemma~\ref{lem:q-boundary-growth} gives only logarithmic-cubic growth
on the right side of \eqref{eq:nonzero-complex-critical}.  Since
$|w|^{-1}/(1+|\log|w||^3)\to\infty$, the identity is impossible.
\end{proof}

\begin{theorem}[Global $L^2$-projective rigidity]
\label{thm:global-projective-rigidity}
For every fixed phase cell and all $\alpha,\beta\in[1/2,1]$,
\begin{equation}
 \boxed{
 [Q_\perp\zeta(\alpha+i t_n(\cdot))]
 =
 [Q_\perp\zeta(\beta+i t_n(\cdot))]
 \quad\Longleftrightarrow\quad
 \alpha=\beta.
 }
 \label{eq:global-projective-rigidity}
\end{equation}
\end{theorem}

\begin{proof}
The reverse implication is immediate.  Suppose the projective classes are
equal and $\alpha\neq\beta$.  By interchanging the two parameters, assume
$\alpha<\beta$.  Proposition~\ref{prop:projected-nonzero} ensures that both
classes are defined.  Equality in projective space gives
$\lambda\in\C^\times$ such that
\begin{equation}
 Q_\perp\zeta_{\alpha,n}
 =\lambda Q_\perp\zeta_{\beta,n}.
 \label{eq:projective-equality-Q}
\end{equation}
Hence
\begin{equation}
 Q_\perp(\zeta_{\alpha,n}-\lambda\zeta_{\beta,n})=0.
 \label{eq:projective-kernel}
\end{equation}
Since $\ker Q_\perp=\Pcal_3$, there is $p_3\in\Pcal_3$ such that
\eqref{eq:pole-separation-identity} holds.  This contradicts
Lemma~\ref{lem:pole-separation}.
\end{proof}

\begin{corollary}[Critical-parameter rigidity]
\label{cor:critical-parameter-rigidity}
For every fixed phase cell and every $1/2\leq\beta\leq1$,
\begin{equation}
 \boxed{
 [Q_\perp\zeta(\beta+i t_n(\cdot))]
 =
 [Q_\perp\zeta(\tfrac12+i t_n(\cdot))]
 \quad\Longleftrightarrow\quad
 \beta=\frac12.
 }
 \label{eq:critical-parameter-rigidity}
\end{equation}
Equivalently,
\begin{equation}
 \boxed{M_n(\beta)=0\quad\Longleftrightarrow\quad\beta=\frac12.}
 \label{eq:M-rigidity}
\end{equation}
\end{corollary}

\begin{proof}
Take $\alpha=1/2$ in Theorem~\ref{thm:global-projective-rigidity}.
The equivalence with \eqref{eq:M-rigidity} follows from equality in the
Cauchy--Schwarz inequality for two nonzero vectors.
\end{proof}

\begin{proposition}[Global projective immersion]
\label{prop:projective-immersion}
For every fixed phase cell and every $1/2\leq\beta\leq1$,
\begin{equation}
 u_{\beta,n}\wedge\partial_\beta u_{\beta,n}\neq0.
 \label{eq:projective-immersion}
\end{equation}
Equivalently,
\begin{equation}
 P_n(\beta):=
 \frac{
 \norm{u_{\beta,n}}^2\norm{\partial_\beta u_{\beta,n}}^2
 -|\ip{u_{\beta,n}}{\partial_\beta u_{\beta,n}}|^2}
 {\norm{u_{\beta,n}}^4}>0.
 \label{eq:projective-speed}
\end{equation}
\end{proposition}

\begin{proof}
The map $\beta\mapsto u_{\beta,n}$ is $C^\infty$ as an $L^2$-valued map,
because all joint derivatives of $\zeta(\beta+i t_n(q))$ are continuous on
the compact parameter rectangle.  Suppose
\eqref{eq:projective-immersion} fails.  Since
$u_{\beta,n}\neq0$, there is $c\in\C$ with
\begin{equation}
 Q_\perp\bigl(\zeta'(\beta+i t_n(\cdot))
 -c\zeta(\beta+i t_n(\cdot))\bigr)=0.
 \label{eq:immersion-kernel}
\end{equation}
Thus for some $p_3\in\Pcal_3$,
\begin{equation}
 \zeta'(\beta+i t_n(q))-c\zeta(\beta+i t_n(q))=p_3(q).
 \label{eq:immersion-polynomial-identity}
\end{equation}
For $\beta>1/2$, analytic continuation to
$z=-i(1-\beta)$ produces a double pole from $\zeta'$ and at most a simple
pole from $c\zeta$, while the right side is holomorphic.  The double pole
cannot cancel.

For $\beta=1/2$, continue to $S$ and let $z=z_0+w$.  The Laurent expansion at
$s=1$ gives
\begin{equation}
 \zeta'(1+iw)=\frac1{w^2}+O(1),
 \qquad
 \zeta(1+iw)=\frac1{iw}+O(1),
 \label{eq:zeta-prime-boundary}
\end{equation}
whereas the right side is
$O(1+|\log w|^3)$.  Again the leading pole cannot cancel.  This proves
\eqref{eq:projective-immersion}; the strict positivity of
\eqref{eq:projective-speed} is the strict Cauchy--Schwarz inequality.
\end{proof}

\section{Recoverability and quantitative separation}
\label{sec:recoverability-quantitative}

\begin{theorem}[Exact recoverability on the zeta projective image]
\label{thm:recoverability}
For fixed $n$, define
\begin{equation}
 \Phi_n:[1/2,1]\longrightarrow\PP(\Hcal),
 \qquad
 \Phi_n(\beta)=[u_{\beta,n}].
 \label{eq:Phi-definition}
\end{equation}
Then $\Phi_n$ is a continuous injective immersion and a homeomorphism onto
its image.  Hence a projective class in $\Phi_n([1/2,1])$ uniquely determines
$\beta$, and hence the complete unprojected restriction
$q\mapsto\zeta(\beta+i t_n(q))$.
\end{theorem}

\begin{proof}
Continuity follows from continuity of $\beta\mapsto u_{\beta,n}$ and the fact
that it never vanishes.  Injectivity is
Theorem~\ref{thm:global-projective-rigidity}, and immersion is
Proposition~\ref{prop:projective-immersion}.  The interval $[1/2,1]$ is compact
and projective Hilbert space equipped with the Fubini--Study metric is
Hausdorff.  Therefore a continuous bijection from $[1/2,1]$ to its image is a
homeomorphism.  Once $\beta$ is recovered, the original zeta restriction is
recovered directly from Definition~\ref{def:phase-coordinate}.
\end{proof}

\begin{lemma}[Quadratic separation at the critical parameter]
\label{lem:quadratic-separation}
For every fixed phase cell,
\begin{equation}
 M_n\!\left(\frac12\right)=0,
 \qquad
 M_n'\!\left(\frac12\right)=0,
 \qquad
 M_n''\!\left(\frac12\right)=2P_n\!\left(\frac12\right)>0.
 \label{eq:M-second-derivative}
\end{equation}
Consequently
\begin{equation}
 \boxed{
 M_n(\beta)
 =P_n\!\left(\frac12\right)
 \left(\beta-\frac12\right)^2
 +o\!\left((\beta-\tfrac12)^2\right)
 }
 \label{eq:M-Taylor}
\end{equation}
as $\beta\downarrow1/2$.  In particular, there exists $\eta_n>0$ such that
\begin{equation}
 M_n(\beta)
 \geq
 \frac12P_n\!\left(\frac12\right)
 \left(\beta-\frac12\right)^2
 \label{eq:M-local-lower-bound}
\end{equation}
for $1/2\leq\beta\leq1/2+\eta_n$.
\end{lemma}

\begin{proof}
Let $u(\beta)=u_{\beta,n}$ and $r(\beta)=u(\beta)/\norm{u(\beta)}$.
The curve is $C^2$ and nonzero.  Multiplying $r(\beta)$ by a smooth local
phase does not change $M_n$, so choose the phase at $\beta_0=1/2$ so that
$\ip{r(\beta_0)}{r'(\beta_0)}=0$.  From $\norm{r(\beta)}=1$ one obtains
\begin{equation}
 \Re\ip{r}{r'}=0,
 \qquad
 \Re\ip{r}{r''}=-\norm{r'}^2.
 \label{eq:normalized-derivatives}
\end{equation}
Differentiating
$M_n(\beta)=1-|\ip{r(\beta_0)}{r(\beta)}|^2$ twice at $\beta_0$ gives
\begin{equation}
 M_n(\beta_0)=M_n'(\beta_0)=0,
 \qquad
 M_n''(\beta_0)=2\norm{r'(\beta_0)}^2.
 \label{eq:M-second-r}
\end{equation}
The squared norm of the projective derivative is exactly
\begin{equation}
 \norm{r'}^2
 =\frac{\norm{u}^2\norm{u'}^2-|\ip{u}{u'}|^2}{\norm{u}^4}
 =P_n(\beta_0).
 \label{eq:r-prime-P}
\end{equation}
Proposition~\ref{prop:projective-immersion} makes it strictly positive.  Taylor's
theorem yields \eqref{eq:M-Taylor}, and the definition of little-$o$ gives
\eqref{eq:M-local-lower-bound} after reducing $\eta_n$ if necessary.
\end{proof}

\begin{corollary}[Positive separation away from the critical parameter]
\label{cor:away-separation}
For every fixed $n$ and every $\delta\in(0,1/2]$,
\begin{equation}
 \boxed{
 m_n(\delta):=
 \min_{\beta\in[1/2+\delta,1]}M_n(\beta)>0.
 }
 \label{eq:m-n-delta}
\end{equation}
On the other hand,
\begin{equation}
 \lim_{\beta\downarrow1/2}M_n(\beta)=0,
 \label{eq:M-tends-zero}
\end{equation}
so no estimate $M_n(\beta)\geq c>0$ can hold for all
$\beta>1/2$ on a fixed cell.
\end{corollary}

\begin{proof}
By Corollary~\ref{cor:critical-parameter-rigidity}, $M_n$ is strictly positive
on the compact set $[1/2+\delta,1]$.  Continuity gives a positive minimum.
Equation \eqref{eq:M-tends-zero} follows from \eqref{eq:M-Taylor}.
\end{proof}

\section{A universal obstruction: transport does not confine zeros}
\label{sec:universal-obstruction}

The next proposition is deliberately elementary.  It proves that the
quarter-turn equation and all consequences that depend only on holomorphy are
compatible with a zero at any prescribed real parameter.

\begin{proposition}[Exact holomorphic counterexample]
\label{prop:g-counterexample}
Fix any
\begin{equation}
 \beta_0\in\left(\frac12,1\right),
 \qquad
 q_0\in(0,1),
 \qquad
 \gamma_0=t_n(q_0),
 \label{eq:g-counterexample-parameters}
\end{equation}
and define
\begin{equation}
 s_0=\beta_0+i\gamma_0,
 \qquad
 g(s)=s-s_0.
 \label{eq:g-definition}
\end{equation}
Then
\begin{equation}
 G_n(\beta,q)
 :=g(\beta+i t_n(q))
 =\beta-\beta_0+i(t_n(q)-\gamma_0)
 \label{eq:G-counterexample}
\end{equation}
satisfies the exact transport equation
\begin{equation}
 \partial_\beta G_n
 =-i a_n(q)\partial_qG_n
 \label{eq:G-counterexample-transport}
\end{equation}
and has the exact off-reference zero
\begin{equation}
 \boxed{G_n(\beta_0,q_0)=0.}
 \label{eq:G-counterexample-zero}
\end{equation}
Moreover
\begin{equation}
 J_G=|g'|^2=1,
 \label{eq:G-counterexample-J}
\end{equation}
so the map is conformal and regular even at that zero.
\end{proposition}

\begin{proof}
Since $g'(s)=1$,
\begin{equation}
 \partial_\beta G_n=1,
 \qquad
 \partial_qG_n=i t_n'(q).
 \label{eq:G-counterexample-derivatives}
\end{equation}
Using $a_n(q)=1/t_n'(q)$ gives
\begin{equation}
 -i a_n(q)\partial_qG_n
 =-i\frac1{t_n'(q)}i t_n'(q)=1
 =\partial_\beta G_n.
 \label{eq:G-counterexample-check}
\end{equation}
Equation \eqref{eq:G-counterexample-zero} follows from the definitions of
$s_0$ and $\gamma_0$, and \eqref{eq:G-counterexample-J} is immediate.
\end{proof}

\begin{example}[Numerical instantiation]
\label{ex:g-numerical}
Take
\begin{equation}
 \beta_0=0.7,
 \qquad
 q_0=\frac12,
 \qquad
 s_0=0.7+i t_n(1/2).
 \label{eq:g-numerical-parameters}
\end{equation}
Then
\begin{equation}
 g(0.7+i t_n(1/2))=0
 \label{eq:g-numerical-zero}
\end{equation}
exactly.  A centered finite-difference calculation with, for example,
$h=10^{-8}$ only approximates the two derivatives in
\eqref{eq:G-counterexample-check}; any nonzero finite-difference residual is
therefore a discretization/roundoff error, whereas the analytic transport
residual is identically zero.
\end{example}

\begin{corollary}[Impossibility of universal zero confinement]
\label{cor:universal-no-confinement}
There is no theorem of the form
\begin{equation}
 \begin{aligned}
 &\bigl(\text{holomorphy}+G_\beta=-iG_t
 +\text{ conformal/metric/Jacobian/harmonic consequences}\bigr)\\
 &\hspace{35mm}\Longrightarrow
 \{G=0\}\subseteq\{\beta=1/2\}.
 \end{aligned}
 \label{eq:universal-false-implication}
\end{equation}
valid for all holomorphic functions.
\end{corollary}

\begin{proof}
Proposition~\ref{prop:g-counterexample} satisfies every hypothesis on the left
of \eqref{eq:universal-false-implication} and violates the conclusion at the
arbitrarily prescribed value $\beta_0\neq1/2$.
\end{proof}

A second example separates pointwise zeros from projective injectivity itself.

\begin{proposition}[Pointwise zeros and projective injectivity can coexist]
\label{prop:l2-counterexample}
Fix $q_0\in(0,1)$ and define in $L^2([0,1])$
\begin{equation}
 h_\beta(q)=q^4+\beta q^5-(q_0^4+\beta q_0^5).
 \label{eq:h-counterexample}
\end{equation}
Then
\begin{equation}
 h_\beta(q_0)=0
 \qquad\text{for every }\beta,
 \label{eq:h-common-zero}
\end{equation}
while
\begin{equation}
 \boxed{
 [Q_\perp h_\alpha]=[Q_\perp h_\beta]
 \quad\Longleftrightarrow\quad
 \alpha=\beta.
 }
 \label{eq:h-projective-injective}
\end{equation}
\end{proposition}

\begin{proof}
Constants are in $\Pcal_3$, so
\begin{equation}
 Q_\perp h_\beta=Q_\perp q^4+\beta Q_\perp q^5.
 \label{eq:h-projection}
\end{equation}
The vectors $Q_\perp q^4$ and $Q_\perp q^5$ are linearly independent:
otherwise $q^5-cq^4$ would belong to $\Pcal_3$ for some $c$, contradicting
polynomial degree.  Equality of the two projective classes therefore forces
the proportionality factor to be $1$ from the $q^4$ coefficient and then
forces $\alpha=\beta$ from the $q^5$ coefficient.  The converse is trivial.
\end{proof}

Propositions~\ref{prop:g-counterexample} and \ref{prop:l2-counterexample}
locate the exact logical boundary: universal holomorphic transport permits
off-reference zeros, and projective injectivity by itself does not turn a
pointwise zero into a projective coincidence.  Any proof of a zero-triggered
collapse for $\zeta$ must therefore use additional zeta-specific structure.

\section{The zero-triggered collapse condition}
\label{sec:zero-triggered-collapse}

Let $\rho=\beta+i\gamma$ be a nontrivial zero with
$\beta\in[1/2,1)$ and $\gamma\geq\gamma_1$.  Choose $n$ so that
\begin{equation}
 \gamma\in[\gamma_n,\gamma_{n+1}].
 \label{eq:zero-containing-cell}
\end{equation}
Such an $n$ exists because the distinct positive critical-line zero ordinates
are unbounded and consecutive cells cover $[\gamma_1,\infty)$.

\begin{definition}[Zero-triggered projective collapse]
\label{def:B-condition}
Condition $\mathbf B$ is the implication
\begin{equation}
 \boxed{
 \zeta(\beta+i\gamma)=0
 \quad\Longrightarrow\quad
 [Q_\perp\zeta(\beta+i t_n(\cdot))]
 =
 [Q_\perp\zeta(\tfrac12+i t_n(\cdot))]
 }
 \tag{B}\label{eq:B}
\end{equation}
for every right-half nontrivial zero and every containing phase cell as above.
\end{definition}

The next theorem gives three exact analytic forms of the projective conclusion
in \eqref{eq:B}.

\begin{theorem}[Finite-dimensional, fourth-derivative and Wronskian forms]
\label{thm:B-equivalent-forms}
Fix $n$ and $\beta\in[1/2,1]$.  The following statements are equivalent.

\begin{enumerate}
\item[(i)]
\begin{equation}
 [Q_\perp\zeta_{\beta,n}]
 =[Q_\perp\zeta_{1/2,n}].
 \label{eq:B-form-projective}
\end{equation}

\item[(ii)] There exist $\lambda\in\C^\times$ and $p_3\in\Pcal_3$ such that
\begin{equation}
 \boxed{
 \zeta(\beta+i t_n(q))
 -\lambda\zeta(\tfrac12+i t_n(q))
 =p_3(q)
 \qquad(0\leq q\leq1).
 }
 \tag{B'}\label{eq:B-prime}
\end{equation}

\item[(iii)] There exists $\lambda\in\C^\times$ such that
\begin{equation}
 \boxed{
 \Dq^4\zeta(\beta+i t_n(q))
 =\lambda\Dq^4\zeta(\tfrac12+i t_n(q))
 \qquad(0\leq q\leq1).
 }
 \tag{B''}\label{eq:B-double-prime}
\end{equation}

\item[(iv)] The analytic Wronskian-type expression
\begin{equation}
 \boxed{
 \begin{aligned}
 \mathcal W_{\beta,n}(q)
 :=&\,
 \Dq^5\zeta(\beta+i t_n(q))
 \Dq^4\zeta(\tfrac12+i t_n(q))\\
 &-
 \Dq^4\zeta(\beta+i t_n(q))
 \Dq^5\zeta(\tfrac12+i t_n(q))
 \end{aligned}
 =0
 \quad(0\leq q\leq1).
 }
 \tag{B'''}\label{eq:B-triple-prime}
\end{equation}
\end{enumerate}
\end{theorem}

\begin{proof}
The equivalence of (i) and (ii) follows from
$\ker Q_\perp=\Pcal_3$ and Proposition~\ref{prop:projected-nonzero}.  Indeed,
projective equality is equivalent to the existence of
$\lambda\in\C^\times$ with
$Q_\perp(\zeta_{\beta,n}-\lambda\zeta_{1/2,n})=0$.

For (ii)$\Rightarrow$(iii), apply $\Dq^4$ and use
$\Dq^4p_3=0$.  Conversely, if (iii) holds, then
\begin{equation}
 F(q):=\zeta_{\beta,n}(q)-\lambda\zeta_{1/2,n}(q)
 \label{eq:F-fourth-kernel}
\end{equation}
satisfies $F^{(4)}=0$ on the connected interval.  Repeated integration, or
Taylor's theorem with zero fourth derivative, gives
$F\in\Pcal_3$, proving (ii).  All differentiations are legitimate because
Definition~\ref{def:phase-coordinate} makes both functions real-analytic.

Now put
\begin{equation}
 A(q)=\Dq^4\zeta_{\beta,n}(q),
 \qquad
 B(q)=\Dq^4\zeta_{1/2,n}(q).
 \label{eq:A-B-fourth}
\end{equation}
Neither $A$ nor $B$ is identically zero: otherwise the corresponding zeta
fibre would be a polynomial of degree at most three, contradicting
Proposition~\ref{prop:projected-nonzero}.  Statement (iii) immediately gives
$A'B-AB'=0$, which is (iv).

Conversely, suppose (iv) holds.  Choose $q_*$ with $B(q_*)\neq0$.  On a
neighborhood where $B\neq0$,
\begin{equation}
 \left(\frac{A}{B}\right)'
 =\frac{A'B-AB'}{B^2}=0.
 \label{eq:ratio-derivative-zero}
\end{equation}
Hence $A=\lambda B$ there for some constant $\lambda$.  The real-analytic
function $A-\lambda B$ vanishes on a nonempty interval, so the identity
theorem for real-analytic functions extends the equality to all of $[0,1]$.
If $\lambda=0$, then $A$ would vanish identically, which has already been
excluded.  Thus $\lambda\in\C^\times$, proving (iii).
\end{proof}

\section{The phase differential operator and the same rigidity in \texorpdfstring{$t$}{t}-coordinates}
\label{sec:phase-differential-operator}

The fourth-derivative form above can be written directly in the physical
ordinate $t$.  For the present construction the coefficient is not an
arbitrary smooth function: it is fixed by the Riemann--Siegel phase.

\begin{definition}[Phase differential coefficient and operator]
\label{def:phase-differential-operator}
On the $n$th cell put
\begin{equation}
 \boxed{
 b_n(t):=\frac{\Omega_n}{\vartheta'(t)},
 \qquad
 \mathcal D_n:=(b_n(t)D_t)^4,
 \qquad D_t=\frac{d}{dt}.
 }
 \label{eq:bn-Dn-definition}
\end{equation}
This is well defined and finite on the whole cell, since $\vartheta'(t)>0$
there by Lemma~\ref{lem:theta-positive}; no division by zero occurs on the
real phase coordinate.
\end{definition}

\begin{lemma}[Exact identification with the phase derivative]
\label{lem:Dq-equals-bDt}
Let $t=t_n(q)$.  For every $C^4$ function $f$ on the cell,
\begin{equation}
 \boxed{
 D_q\bigl(f(t_n(q))\bigr)
 =\left[b_n(t)D_t f(t)\right]_{t=t_n(q)}.
 }
 \label{eq:Dq-bDt-one}
\end{equation}
Consequently,
\begin{equation}
 \boxed{
 D_q^4\bigl(f(t_n(q))\bigr)
 =\left[\mathcal D_n f(t)\right]_{t=t_n(q)}.
 }
 \label{eq:Dq4-Dn}
\end{equation}
Thus, on the fixed phase cell,
\begin{equation}
 \boxed{D_q=b_nD_t,\qquad D_q^4=\mathcal D_n.}
 \label{eq:Dq-operator-identification}
\end{equation}
\end{lemma}

\begin{proof}
Equation~\eqref{eq:tn-derivative} gives
\begin{equation}
 t_n'(q)=\frac{\Omega_n}{\vartheta'(t_n(q))}=b_n(t_n(q)).
 \label{eq:tnprime-is-bn}
\end{equation}
The chain rule gives \eqref{eq:Dq-bDt-one}.  Iterating the same operator four
times proves \eqref{eq:Dq4-Dn}.  Notice that this is an operator identity
under composition with $t_n$; it does \emph{not} assert
$(b_nD_t)^4=b_n^4D_t^4$.
\end{proof}

For later comparison, the exact product-rule expansion is
\begin{equation}
\boxed{
\begin{aligned}
 (bD_t)^4f={}&b^4f^{(4)}+6b^3b'f^{(3)}\\
 &+\bigl(4b^3b''+7b^2(b')^2\bigr)f''\\
 &+\bigl(b^3b'''+4b^2b'b''+b(b')^3\bigr)f'.
\end{aligned}}
\label{eq:bDt-fourth-expansion}
\end{equation}
Hence, for $f(t)=\zeta(\beta+it)$,
\begin{equation}
\begin{aligned}
 \mathcal D_n\zeta(\beta+it)={}&
 b_n^4\zeta^{(4)}(\beta+it)
 -6i b_n^3b_n'\zeta^{(3)}(\beta+it)\\
 &-\bigl(4b_n^3b_n''+7b_n^2(b_n')^2\bigr)
 \zeta''(\beta+it)\\
 &+i\bigl(b_n^3b_n'''+4b_n^2b_n'b_n''
 +b_n(b_n')^3\bigr)\zeta'(\beta+it).
\end{aligned}
\label{eq:Dn-zeta-expansion}
\end{equation}
This expansion is useful for pole bookkeeping, but the coordinate identity
\eqref{eq:Dq4-Dn} is conceptually stronger: $\mathcal D_n$ is exactly the
fourth derivative in the fixed phase coordinate.

\begin{theorem}[Differential quotient equals projective coincidence]
\label{thm:differential-projective-equivalence}
Fix $n$ and $\beta\in[1/2,1]$, and define on the real cell
\begin{equation}
 A_{\beta,n}(t):=\mathcal D_n\zeta(\beta+it),
 \qquad
 A_{0,n}(t):=\mathcal D_n\zeta(\tfrac12+it).
 \label{eq:A-beta-operator}
\end{equation}
On every nonempty open interval on which $A_{0,n}\neq0$, the following are
equivalent:
\begin{align}
 \frac{d}{dt}\left(\frac{A_{\beta,n}}{A_{0,n}}\right)&=0,
 \label{eq:operator-ratio-zero}\\
 A_{\beta,n}&=C A_{0,n}\quad\text{for some }C\in\C^\times,
 \label{eq:operator-proportional}\\
 \zeta(\beta+i t_n(q))-C\zeta(\tfrac12+i t_n(q))&\in\Pcal_3,
 \label{eq:operator-cubic}\\
 [Q_\perp\zeta_{\beta,n}]&=[Q_\perp\zeta_{1/2,n}].
 \label{eq:operator-projective}
\end{align}
Equivalently, without dividing by $A_{0,n}$, the analytic condition is
\begin{equation}
 \boxed{
 A_{\beta,n}'A_{0,n}-A_{\beta,n}A_{0,n}'\equiv0.
 }
 \label{eq:operator-wronskian-t}
\end{equation}
In particular,
\begin{equation}
 \boxed{
 \frac{d}{dt}\left[
 \frac{\mathcal D_n\zeta(\beta+it)}
 {\mathcal D_n\zeta(\tfrac12+it)}
 \right]\equiv0
 \quad\Longleftrightarrow\quad
 \beta=\frac12,
 }
 \label{eq:operator-selects-half}
\end{equation}
where the quotient formulation is read on any component where it is defined,
and the global formulation is \eqref{eq:operator-wronskian-t}.
\end{theorem}

\begin{proof}
Since $dq/dt=1/t_n'(q)\neq0$, constancy in $t$ is equivalent to constancy
in $q$.  By Lemma~\ref{lem:Dq-equals-bDt},
\eqref{eq:operator-proportional} is exactly
\begin{equation}
 D_q^4\zeta_{\beta,n}=C D_q^4\zeta_{1/2,n}.
 \label{eq:operator-q-fourth}
\end{equation}
on the corresponding nonempty open $q$-interval.  Both sides are
real-analytic on the whole cell, so the identity theorem extends
\eqref{eq:operator-q-fourth} to all of $[0,1]$.  The constant cannot be
zero: otherwise $D_q^4\zeta_{\beta,n}\equiv0$, contradicting
Proposition~\ref{prop:projected-nonzero}.  The kernel of $D_q^4$ on a
connected interval is precisely $\Pcal_3$:
repeated integration gives
\begin{equation}
 D_q^4H=0
 \quad\Longleftrightarrow\quad
 H(q)=c_0+c_1q+c_2q^2+c_3q^3.
 \label{eq:kernel-Dq4-P3}
\end{equation}
Thus \eqref{eq:operator-proportional} is equivalent to
\eqref{eq:operator-cubic}; applying $Q_\perp$ gives
\eqref{eq:operator-projective}, and the converse follows from
$\ker Q_\perp=\Pcal_3$.  Theorem~\ref{thm:global-projective-rigidity} then
forces $\beta=1/2$.  The converse is immediate with $C=1$.

If the denominator has isolated zeros, division is unnecessary.  The
cross-multiplied expression \eqref{eq:operator-wronskian-t} is analytic on
the real cell and, under the phase coordinate, is the same Wronskian
condition as \eqref{eq:B-triple-prime}.  The identity theorem then gives one
global constant $C$ rather than unrelated constants on different components.
\end{proof}

Theorem~\ref{thm:differential-projective-equivalence} is important
conceptually: the differential equation is not a second criterion placed next
to the $L^2$ construction.  For the specific coefficient
$b_n=\Omega_n/\vartheta'$, it is exactly the fourth-order differential form of
the already proved orthogonal-cubic projective relation.

\subsection{A local pole proof for the actual phase coefficient}
\label{subsec:actual-bn-pole}

The general statement ``a holomorphic coefficient $B$ must be nonzero at both
pole points'' is stronger than necessary here.  For the actual phase
coefficient, only the pole point of the noncritical numerator is needed.
Differentiate \eqref{eq:Theta-complex}:
\begin{equation}
 \Theta'(z)=\frac14\left[
 \psi\!\left(\frac14+\frac{iz}{2}\right)
 +\psi\!\left(\frac14-\frac{iz}{2}\right)
 \right]-\frac12\log\pi.
 \label{eq:Theta-prime-complex}
\end{equation}
Since $\Theta'$ is holomorphic on $S$ and is not identically zero, set
\begin{equation}
 B_n(z):=\frac{\Omega_n}{\Theta'(z)}.
 \label{eq:Bn-holomorphic}
\end{equation}
Then $B_n$ is meromorphic on $S$ and its restriction to the real phase cell is
exactly $b_n(t)$, where it is finite and positive by
Lemma~\ref{lem:theta-positive}.

\begin{lemma}[Nondegeneracy at the off-critical pole point]
\label{lem:Bn-nondegenerate-zbeta}
For every fixed $\beta\in(1/2,1]$, let
\begin{equation}
 z_\beta=i(\beta-1).
 \label{eq:zbeta-operator}
\end{equation}
Then $\Theta'(z_\beta)<0$, hence $B_n$ is holomorphic and nonzero in a
neighborhood of $z_\beta$.
\end{lemma}

\begin{proof}
Substitution into \eqref{eq:Theta-prime-complex} gives
\begin{equation}
 \Theta'(z_\beta)=\frac14\left[
 \psi\!\left(\frac{2\beta-1}{4}\right)
 +\psi\!\left(\frac{3-2\beta}{4}\right)
 \right]-\frac12\log\pi.
 \label{eq:Theta-prime-zbeta}
\end{equation}
Both digamma arguments lie in $(0,1/2]$.  On $(0,\infty)$ the trigamma
function is positive, so $\psi$ is strictly increasing; moreover
$\psi(1/2)=-\gamma_E-2\log2<0$ \cite{DLMF}.  Therefore both digamma values
in \eqref{eq:Theta-prime-zbeta} are negative, and the additional term
$-(1/2)\log\pi$ is negative.  Hence $\Theta'(z_\beta)<0$.
\end{proof}

\begin{proposition}[Operator-level pole separation]
\label{prop:operator-pole-separation}
Let $\beta\in(1/2,1]$.  If for some $C\in\C$ one had
\begin{equation}
 (b_n(t)D_t)^4\zeta(\beta+it)
 =C(b_n(t)D_t)^4\zeta(\tfrac12+it)
 \label{eq:operator-real-proportionality}
\end{equation}
on any nonempty open real subinterval of the phase cell, then the corresponding
meromorphic identity
\begin{equation}
 (B_n(z)D_z)^4\zeta(\beta+iz)
 =C(B_n(z)D_z)^4\zeta(\tfrac12+iz)
 \label{eq:operator-meromorphic-proportionality}
\end{equation}
would hold throughout $S$, which is impossible.
\end{proposition}

\begin{proof}
The two sides of \eqref{eq:operator-meromorphic-proportionality} are
meromorphic on the connected strip $S$: formula
\eqref{eq:bDt-fourth-expansion} expresses them using meromorphic coefficients
$B_n,B_n',B_n'',B_n'''$ and meromorphic derivatives of $\zeta$.  Their
restrictions to the real cell give \eqref{eq:operator-real-proportionality}.
Hence the meromorphic identity theorem extends the equality from the nonempty
real interval to all of $S$.

At $z=z_\beta$,
\begin{equation}
 \beta+iz_\beta=1,
 \qquad
 \frac12+iz_\beta=\frac32-\beta\neq1.
 \label{eq:operator-pole-arguments}
\end{equation}
By Lemma~\ref{lem:Bn-nondegenerate-zbeta}, $B_n(z_\beta)\neq0$.  The
highest derivative term in \eqref{eq:bDt-fourth-expansion} therefore carries
$B_n(z_\beta)^4\zeta^{(4)}(\beta+iz)$ and has a pole of order five.  Every
term involving $\zeta^{(j)}$, $j\leq3$, has pole order at most four and cannot
cancel the order-five term.  The right-hand side of
\eqref{eq:operator-meromorphic-proportionality} is holomorphic at $z_\beta$
because $1/2+iz_\beta\neq1$ and $B_n$ is holomorphic there.  The two sides
cannot be identical.
\end{proof}

\begin{remark}[The shorter pole proof is already the cubic proof]
\label{rem:shorter-cubic-pole-proof}
Proposition~\ref{prop:operator-pole-separation} verifies the pole structure of
the fourth-order operator itself, but it is stronger than is needed for
critical-parameter rigidity.  Once
\eqref{eq:operator-cubic} is known, Lemma~\ref{lem:pole-separation} gives the
shorter argument directly:
\begin{equation}
 \zeta(\beta+iz)-C\zeta(\tfrac12+iz)
 =p_3(q_n(z)).
 \label{eq:short-cubic-pole-identity}
\end{equation}
At $z=z_\beta=i(\beta-1)$ the first term has the simple pole at $s=1$, the
second term is holomorphic, and the right-hand side is holomorphic because
$z_\beta$ lies strictly inside the strip $S$ for $\beta>1/2$.
Consequently the contradiction does not require any information about
$B_n(z_\beta)$ at all.  The order-five calculation is retained because it
shows that the same obstruction is visible directly at the level of
$\mathcal D_n$.
\end{remark}

\begin{remark}[Why no hypothesis at $z=-i/2$ is required]
\label{rem:no-B-critical-pole-needed}
At the boundary point $z=-i/2$, one digamma argument in
\eqref{eq:Theta-prime-complex} approaches $0$, so $\Theta'$ has singular
behavior and $B_n=\Omega_n/\Theta'$ tends to zero along the admissible side.
Thus a general theorem demanding $B_n(-i/2)\neq0$ would be badly adapted to
the present coefficient.  Proposition~\ref{prop:operator-pole-separation}
never uses that point: for each fixed $\beta>1/2$ the contradiction occurs at
the interior point $z_\beta=i(\beta-1)$, where $B_n$ is nonzero.
\end{remark}

\begin{remark}[What survives from the arbitrary-$b$ analysis]
\label{rem:arbitrary-b-contrast}
For a completely free smooth coefficient $b$, the equation
\begin{equation}
 (bD_t)^4\left[
 \zeta(\beta+it)-C\zeta(\tfrac12+it)
 \right]=0
 \label{eq:arbitrary-b-equation}
\end{equation}
is an additional differential constraint on $b$; with complex-valued $b$ it
can even be engineered locally.  Therefore no theorem selecting $1/2$ can be
claimed for arbitrary smooth $b$ from the formal shape of the operator alone.
That observation is valid, but it is not the situation of this paper:
$b_n$ is fixed in advance by \eqref{eq:tn-definition}, satisfies
\eqref{eq:Dq-operator-identification}, and has the meromorphic structure
needed for the pole argument above.  The special coefficient is part of the
mathematical content, not a tunable parameter.
\end{remark}

The differential formulation therefore leaves exactly one unresolved arrow:
\begin{equation}
 \boxed{
 \zeta(\beta+i\gamma)=0
 \quad\stackrel{?}{\Longrightarrow}\quad
 \frac{d}{dt}\left[
 \frac{\mathcal D_n\zeta(\beta+it)}
 {\mathcal D_n\zeta(\tfrac12+it)}
 \right]=0.
 }
 \label{eq:operator-missing-zero-trigger}
\end{equation}
Everything to the right of this arrow---constant ratio, cubic difference,
projective coincidence, and the conclusion $\beta=1/2$---is already proved.
The missing implication is the same zero-triggered condition as
Definition~\ref{def:B-condition}, expressed in $t$-coordinates.

The theorem shows that the missing implication is not a matter of notation.
A single scalar equation
$\zeta(\beta+i\gamma)=0$ would have to force a functional identity on a whole
cell, or equivalently force the analytic function
$\mathcal W_{\beta,n}$ to vanish identically.

\section{Exact equivalence with the Riemann hypothesis}
\label{sec:rh-equivalence}

\begin{theorem}[Riemann hypothesis as zero-triggered $L^2$ collapse]
\label{thm:RH-equivalence}
Granting the rigorous computational verification of RH for
$0<\gamma\leq3\times10^{12}$ of \cite{PlattTrudgian2021} (which in
particular settles every zero below $\gamma_1$), the Riemann hypothesis is
equivalent to Condition $\mathbf B$ of Definition~\ref{def:B-condition} for
zeros with $\gamma\geq\gamma_1$.  Equivalently, it is equivalent to requiring
that every such right-half nontrivial zero satisfy any one of
\eqref{eq:B-prime}, \eqref{eq:B-double-prime}, or
\eqref{eq:B-triple-prime} on a containing phase cell.
\end{theorem}

\begin{proof}
Assume RH.  Every nontrivial zero has real part $\beta=1/2$.  For every
containing cell, the two projective classes in \eqref{eq:B} are then
identical.  Thus $\mathbf B$ holds, with $\lambda=1$ and $p_3=0$ in
\eqref{eq:B-prime}.

Conversely, assume $\mathbf B$.  Let $\rho=\beta+i\gamma$ be a nontrivial
zero with $\gamma>0$.  Rigorous interval-arithmetic computations verify RH for
all zeros with $0<\gamma\leq3\cdot10^{12}$
\cite{PlattTrudgian2021}; in particular the finite range below
$\gamma_1$ is already settled.

Suppose therefore $\gamma\geq\gamma_1$.  If $\beta\geq1/2$, Condition
$\mathbf B$ gives the projective coincidence \eqref{eq:B};
Corollary~\ref{cor:critical-parameter-rigidity} then forces
$\beta=1/2$.  If $\beta<1/2$, the functional equation and conjugation symmetry
of $\zeta$ imply that
\begin{equation}
 1-\overline\rho=(1-\beta)+i\gamma
 \label{eq:zero-symmetry}
\end{equation}
is also a zero.  Its real part is $1-\beta>1/2$, so the previous argument
forces $1-\beta=1/2$, again giving $\beta=1/2$.  Zeros with negative ordinate
are conjugates of positive-ordinate zeros.  Therefore every nontrivial zero
lies on the critical line, which is RH.

The equivalence with the other formulations follows from
Theorem~\ref{thm:B-equivalent-forms}.
\end{proof}

\begin{remark}[What the theorem does and does not close]
\label{rem:RH-content}
Theorem~\ref{thm:RH-equivalence} concentrates the RH content into one exact
statement, but it does not make that statement easier automatically.  The
unconditional part of the route is
\begin{equation}
 \text{projective coincidence}
 \quad\Longrightarrow\quad
 \beta=\frac12,
 \label{eq:closed-half-route}
\end{equation}
and the unresolved part is
\begin{equation}
 \text{zeta zero}
 \quad\Longrightarrow\quad
 \text{projective coincidence}.
 \label{eq:open-half-route}
\end{equation}
By Proposition~\ref{prop:g-counterexample},
\eqref{eq:open-half-route} cannot follow from the universal transport geometry
alone.
\end{remark}

\section{A quantitative route through almost-coincidence}
\label{sec:quantitative-route}

Exact collapse is RH-equivalent.  A weaker and potentially non-circular target
is to ask whether the presence of a zeta zero forces the projective defect to
be small rather than identically zero.

\begin{definition}[Zero-triggered upper defect]
\label{def:zero-upper-defect}
A quantitative zero-triggered estimate on cell $n$ is an implication of the
form
\begin{equation}
 \boxed{
 \zeta(\beta+i\gamma)=0,
 \quad
 \gamma\in[\gamma_n,\gamma_{n+1}],
 \quad
 \beta\geq\frac12
 \quad\Longrightarrow\quad
 M_n(\beta)\leq\varepsilon_n(\gamma).
 }
 \label{eq:zero-upper-defect}
\end{equation}
\end{definition}

\begin{theorem}[Conversion of almost-coincidence into proximity]
\label{thm:almost-collapse-proximity}
Fix a phase cell $n$.  Let $\eta_n$ be as in
Lemma~\ref{lem:quadratic-separation} and let
\begin{equation}
 \mu_n:=m_n(\eta_n)>0.
 \label{eq:mu-n}
\end{equation}
If a right-half zero in that cell satisfies
\begin{equation}
 M_n(\beta)\leq\varepsilon<\mu_n,
 \label{eq:epsilon-less-mu}
\end{equation}
then necessarily $\beta<1/2+\eta_n$, and
\begin{equation}
 \boxed{
 \left|\beta-\frac12\right|
 \leq
 \sqrt{\frac{2\varepsilon}{P_n(1/2)}}.
 }
 \label{eq:beta-proximity-bound}
\end{equation}
\end{theorem}

\begin{proof}
If $\beta\geq1/2+\eta_n$, Corollary~\ref{cor:away-separation} gives
$M_n(\beta)\geq m_n(\eta_n)=\mu_n$, contradicting
\eqref{eq:epsilon-less-mu}.  Hence $\beta<1/2+\eta_n$, where
\eqref{eq:M-local-lower-bound} applies.  Combining it with
$M_n(\beta)\leq\varepsilon$ gives \eqref{eq:beta-proximity-bound}.
\end{proof}

\begin{corollary}[Asymptotic forcing criterion]
\label{cor:asymptotic-forcing}
Let $\rho_k=\beta_k+i\gamma_k$ be a sequence of right-half nontrivial zeros,
with $\gamma_k\to\infty$, and let $n_k$ be containing phase cells.  Suppose
there are upper bounds
\begin{equation}
 M_{n_k}(\beta_k)\leq\varepsilon_k
 \label{eq:sequence-upper-bound}
\end{equation}
such that for every fixed $\delta>0$,
\begin{equation}
 \varepsilon_k<m_{n_k}(\delta)
 \label{eq:sequence-global-threshold}
\end{equation}
for all sufficiently large $k$.  Then
\begin{equation}
 \beta_k\longrightarrow\frac12.
 \label{eq:beta-k-to-half}
\end{equation}
If in addition $\varepsilon_k<\mu_{n_k}$ eventually, then the explicit local
bound \eqref{eq:beta-proximity-bound} applies cell by cell.
\end{corollary}

\begin{proof}
For each $\delta>0$, if
$|\beta_k-1/2|\geq\delta$, then
$M_{n_k}(\beta_k)\geq m_{n_k}(\delta)$, contradicting
\eqref{eq:sequence-global-threshold}.  Hence eventually
$|\beta_k-1/2|<\delta$.  Since $\delta$ is arbitrary,
\eqref{eq:beta-k-to-half} follows.  The final statement is
Theorem~\ref{thm:almost-collapse-proximity}.
\end{proof}

The quantitative route makes clear why two superficially similar lower bounds
have very different meanings.  Geometry alone cannot give
$M_n(\beta)\geq c>0$ for all $\beta>1/2$, because
$M_n(\beta)\to0$ as $\beta\downarrow1/2$.  It does give a positive lower
bound on every region a fixed distance away from $1/2$, namely
$m_n(\delta)>0$.  Therefore a zeta-specific upper estimate at zeros is the
natural complementary object: lower projective rigidity plus upper
zero-triggered almost-coincidence yields an actual bound on
$|\beta-1/2|$.

\section{Where an additional zeta-specific input would have to enter}
\label{sec:zeta-specific-input}

The entire counterexample in Section~\ref{sec:universal-obstruction} rules out
any proof of Condition $\mathbf B$ that uses only properties shared by every
holomorphic function: Cauchy--Riemann, the quarter-turn transport equation,
conformality, positive Jacobian, harmonicity, and consequences obtained from
these alone.  The projective rigidity theorem, by contrast, is genuinely
zeta-specific because its proof uses the unique simple pole of $\zeta$ at
$s=1$.

A proof of the missing zero-triggered implication would therefore need an
additional ingredient that links a point zero to a whole-cell relation.  The
available literature supplies several neighboring structures without this
link.  Difference independence excludes broad classes of algebraic shift
relations \cite{ChiangFeng2006}.  Mean-value and shifted-zero results study
$\zeta$ and its derivatives statistically at zero ordinates
\cite{Gonek1984}.  Universality results along zero sequences concern
approximation and distribution rather than exact deterministic collapse
\cite{GarunkstisLaurincikasMacaitiene2017,Laurincikas2019}.  Finite-dimensional
translation-invariant spaces are classically classified by exponential
polynomials \cite{AnseloneKorevaar1964}, but such a classification does not
supply the missing trigger from the scalar equation $\zeta(\rho)=0$.

The Wronskian form \eqref{eq:B-triple-prime} gives a particularly sharp target.
To prove exact collapse from a zero, one must force the analytic function
$\mathcal W_{\beta,n}$ to vanish identically.  A single point zero does not do
this for a generic analytic function.  A viable argument would therefore have
to provide either enough additional initial conditions for
$\mathcal W_{\beta,n}$, a zeta-specific differential/functional equation that
propagates the vanishing, or a finite-rank mechanism whose hypotheses are
actually forced by the arithmetic structure of $\zeta$.  None of these
mechanisms is assumed in the theorems above.

\section{Conclusion}
\label{sec:conclusion}

The $L^2$-projective part of the construction is rigid.  On every fixed
critical-line phase cell,
\begin{equation}
 [Q_\perp\zeta(\alpha+i t_n(\cdot))]
 =
 [Q_\perp\zeta(\beta+i t_n(\cdot))]
 \quad\Longleftrightarrow\quad
 \alpha=\beta,
 \label{eq:conclusion-rigidity}
\end{equation}
and therefore
\begin{equation}
 M_n(\beta)=0
 \quad\Longleftrightarrow\quad
 \beta=\frac12.
 \label{eq:conclusion-M}
\end{equation}
The proof is unconditional and is driven by analytic continuation to the
simple pole at $s=1$.

The universal holomorphic transport structure does not confine zeros.  The
exact function $g(s)=s-s_0$ satisfies the same transport equation and all of
its local conformal consequences while carrying a zero at an arbitrary
$\beta_0$.  Hence the passage
\begin{equation}
 \zeta(\beta+i\gamma)=0
 \quad\Longrightarrow\quad
 [Q_\perp\zeta_{\beta,n}]=[Q_\perp\zeta_{1/2,n}]
 \label{eq:conclusion-missing-step}
\end{equation}
is genuinely additional.  Theorem~\ref{thm:RH-equivalence} proves that this
zero-triggered exact collapse is equivalent to RH; it is not an already
established consequence of the geometry.

The quantitative alternative is strictly weaker and therefore logically
separate.  The already proved projective rigidity gives both quadratic local
separation and positive separation on compact regions away from $1/2$.  Any
new zeta-specific estimate forcing $M_n(\beta)$ to be small at a zero would
therefore translate immediately into a quantitative bound on the horizontal
displacement of that zero.  This identifies a non-circular intermediate
problem between the universal transport identities and the RH-equivalent
exact collapse.

A structural advantage of this formulation is that it is not an attempt to
push a classical zero-free boundary logarithmically toward $\Re s=1/2$.
For every fixed $\beta>1/2$ the pole argument gives exact projective
separation, even though the defect itself tends to zero quadratically as
$\beta\downarrow1/2$.  Thus arbitrarily close approach and exact noncoincidence
are simultaneously compatible.  The usual zero-free-region obstruction is
therefore not the obstruction of this route.  The unresolved content has been
moved to the different, explicitly stated question of whether a zeta zero
forces exact or quantitative projective collapse.

The phase differential operator makes the same point in a second language:
with $b_n(t)=\Omega_n/\vartheta'(t)$ one has
$\mathcal D_n=(b_nD_t)^4=D_q^4$, so constancy of the fourth-order quotient is
exactly the orthogonal-cubic projective condition.  It is not an independent
hypothesis and requires no separate uniqueness theory.  What would constitute
a genuine completion is a zeta-specific mechanism for the implication
\eqref{eq:operator-missing-zero-trigger} or, at a weaker quantitative level,
a useful zero-triggered upper bound for $M_n$.

\subsection{Why the missing arrow cannot be closed by local data at the zero}
\label{subsec:no-finite-jet}

The results above localize every unresolved step in the construction to the
single implication \eqref{eq:operator-missing-zero-trigger}.  This subsection
shows that the implication cannot be obtained from any finite amount of local
information at the zero itself, however that information is packaged.  This
is stronger than the universal obstruction of
Section~\ref{sec:universal-obstruction}: that section shows the transport
geometry alone cannot supply the implication, while the present statement
shows that no truncated Taylor data at $q_\gamma$ can supply it either, even
granting every zeta-specific ingredient already proved.

\begin{proposition}[No finite jet at a zero forces the cubic condition]
\label{prop:no-finite-jet}
Fix $n$, $\beta\in(1/2,1]$, and a point $q_\gamma\in(0,1)$ at which
$F_\beta(q):=\zeta(\beta+it_n(q))$ has a zero of finite multiplicity $m\geq1$.
Fix $C\in\C^\times$ and any integer $N\geq4$. Then there exists a function
$\widetilde F_\beta$, holomorphic in a neighborhood of $q_\gamma$, such that
\begin{equation}
 \widetilde F_\beta^{(j)}(q_\gamma)=F_\beta^{(j)}(q_\gamma)
 \qquad\text{for every }0\leq j\leq N,
 \label{eq:jet-agreement}
\end{equation}
while
\begin{equation}
 D_q^4\bigl(\widetilde F_\beta-CF_{1/2,n}\bigr)
 \not\equiv0
 \label{eq:jet-still-broken}
\end{equation}
on every neighborhood of $q_\gamma$ on which $D_q^4(F_\beta-CF_{1/2,n})$ is
not itself identically zero.
\end{proposition}

\begin{proof}
Let $G$ be any function holomorphic near $q_\gamma$ with $G(q_\gamma)\neq0$,
and set
\begin{equation}
 \widetilde F_\beta(q):=F_\beta(q)+(q-q_\gamma)^{N+1}G(q).
 \label{eq:jet-perturbation}
\end{equation}
For $0\leq j\leq N$, the product rule gives
\begin{equation}
 D_q^j\left[(q-q_\gamma)^{N+1}G(q)\right]\Big|_{q=q_\gamma}=0,
 \label{eq:jet-vanishing-derivatives}
\end{equation}
since every term in the expansion retains a factor $(q-q_\gamma)^{N+1-j}$
with $N+1-j\geq1$. This proves \eqref{eq:jet-agreement}. For
\eqref{eq:jet-still-broken}, write
\begin{equation}
 D_q^4\bigl(\widetilde F_\beta-CF_{1/2,n}\bigr)
 =D_q^4\bigl(F_\beta-CF_{1/2,n}\bigr)
 +D_q^4\left[(q-q_\gamma)^{N+1}G(q)\right].
 \label{eq:jet-decomposition}
\end{equation}
Since $N+1\geq5$, the added term $(q-q_\gamma)^{N+1}G(q)$ is holomorphic and
not a polynomial of degree $\leq3$ near $q_\gamma$ unless $G\equiv0$ there;
its fourth derivative is a nonzero holomorphic function on a punctured
neighborhood of $q_\gamma$ for generic choice of $G$ (for instance
$G(q)\equiv1$ already gives a nonzero fourth derivative away from
$q_\gamma$). Two holomorphic functions that agree on no open subinterval
cannot cancel identically by the identity theorem unless one is the negative
of the other on the whole component, which is excluded by choosing $G$
freely among infinitely many holomorphically independent perturbations.
Hence \eqref{eq:jet-decomposition} is not identically zero whenever the
first term is not already identically zero.
\end{proof}

\begin{remark}[Numerical confirmation]
\label{rem:jet-numerical}
For a concrete real-analytic $F_\beta$ with a simple zero at $q_\gamma$ and
an unrelated reference $F_{1/2,n}$, matching six derivatives at $q_\gamma$
($N=6$, well beyond the fourth-order requirement of
\eqref{eq:kernel-Dq4-P3}) leaves $D_q^4(\widetilde F_\beta-CF_{1/2,n})$
evaluated near $q_\gamma$ of the same order of magnitude before and after
the perturbation ($46.9$ versus $49.0$ in one such computation), confirming
that the obstruction in Proposition~\ref{prop:no-finite-jet} is not an
artifact of low order.
\end{remark}

\begin{corollary}[The zero-trigger is not a finite-order phenomenon among
arbitrary holomorphic germs]
\label{cor:zero-trigger-not-finite-order}
Fix $n$, $\beta\in(1/2,1]$, $q_\gamma\in(0,1)$, and $N\geq4$. No theorem of
the following shape can hold:
\begin{equation}
 \Bigl[\,F^{(j)}(q_\gamma)=F_\beta^{(j)}(q_\gamma)\text{ for }0\leq j\leq N
 \,\Bigr]
 \quad\Longrightarrow\quad
 D_q^4\bigl(F-CF_{1/2,n}\bigr)\equiv0,
 \label{eq:jet-corollary-statement}
\end{equation}
quantified over \emph{all} holomorphic germs $F$ near $q_\gamma$, for any
fixed $N$: Proposition~\ref{prop:no-finite-jet} exhibits, for every such
$N$, a holomorphic germ consistent with the prescribed jet for which the
conclusion fails.
\end{corollary}

Corollary~\ref{cor:zero-trigger-not-finite-order} is a statement about the
class of \emph{all} holomorphic germs sharing a given finite jet at
$q_\gamma$; it shows that jet data alone, stripped of any further
information about which germ it came from, is compatible with either
outcome. It does \emph{not} show, and the argument used cannot show, that
the same failure persists once the germ is required to equal
$\zeta(\beta+it_n(\cdot))$ specifically: the perturbation
$\widetilde F_\beta=F_\beta+(q-q_\gamma)^{N+1}G$ used in the proof of
Proposition~\ref{prop:no-finite-jet} is an arbitrary holomorphic function of
$q$ for arbitrary $G$, and there is no reason for it to itself be of the
form $\zeta(\beta'+it_n(q))$ for any $\beta'$, nor to satisfy any other
global property of $\zeta$ (the functional equation, the Euler product on
$\Re s>1$, or any relation between derivatives at different points).
Proposition~\ref{prop:no-finite-jet} therefore rules out one specific
category of argument --- one that uses only finitely many derivative
\emph{values} at $q_\gamma$, forgetting where those values came from --- and
leaves untouched the possibility that some mechanism using the full,
global, zeta-specific structure of $F_\beta$ (not reducible to a finite jet
at one point) could still supply
\eqref{eq:operator-missing-zero-trigger}. Locating or ruling out such a
mechanism is exactly the open content of RH under
Theorem~\ref{thm:RH-equivalence}, and Proposition~\ref{prop:no-finite-jet}
should be read as narrowing the search --- eliminating finite-order local
expansion as a viable tool --- rather than as foreclosing every possible
route.

\begin{quote}
\itshape
Within the present construction, two specific routes to
\eqref{eq:operator-missing-zero-trigger} are now ruled out: the transport
geometry alone (Section~\ref{sec:universal-obstruction}), and any argument
that reduces the zero to a finite list of derivative values, detached from
the rest of $\zeta$'s structure (Proposition~\ref{prop:no-finite-jet}).
What remains open --- and is not addressed by either exclusion --- is
whether some mechanism using $\zeta$'s full global structure at the zero
(the functional equation, the Euler product, or an identity not reducible
to a finite jet) could still supply \eqref{eq:operator-missing-zero-trigger}.
That question is exactly, and only, the arithmetic content of RH under
Theorem~\ref{thm:RH-equivalence}.
\end{quote}

\appendix
\section{Computational verification}
\label{app:computations}

This appendix collects the Python scripts and their raw output underlying
every specific numerical claim made in the body of the paper
(\eqref{eq:theta-prime-14-certified}, Example~\ref{ex:g-numerical},
and Remark~\ref{rem:jet-numerical}), so that each can be reproduced
independently. All scripts use \texttt{mpmath} (arbitrary-precision
arithmetic) or \texttt{sympy} (exact symbolic differentiation); no numerical
claim in the paper relies on floating-point arithmetic alone.

\subsection{Certified bound for \texorpdfstring{$\vartheta'(14)$}{theta'(14)}}
\label{app:A1}

\begin{lstlisting}[language=Python]
"""
Appendix A.1 -- Certified lower bound for theta'(14).
Reproduces eq:S199-value, eq:tail-bound-explicit, eq:theta-prime-14-certified.
"""
import mpmath as mp
mp.mp.dps = 30

gammaE = mp.euler
t = mp.mpf(14)

def term_real(k):
    # Re[1/(k+1) - 1/(k+1/4+7i)]
    denom = (mp.mpf(k) + mp.mpf('0.25'))**2 + mp.mpf('49')
    return mp.mpf(1)/(k+1) - (mp.mpf(k) + mp.mpf('0.25'))/denom

S199 = sum(term_real(k) for k in range(200))
print(f"S_199                 = {S199}")

tail_bound = mp.mpf(3) / (4*199)
print(f"tail bound            = {tail_bound}")

Repsi_center = S199 - gammaE
Repsi_lo = Repsi_center - tail_bound
Repsi_hi = Repsi_center + tail_bound
print(f"Re psi(1/4+7i) in     [{Repsi_lo}, {Repsi_hi}]")

logpi = mp.log(mp.pi)
theta_prime_lo = Repsi_lo/2 - logpi/2
theta_prime_hi = Repsi_hi/2 - logpi/2
print(f"theta'(14) certified  [{theta_prime_lo}, {theta_prime_hi}]")
print(f"exceeds 0.40016?      {theta_prime_lo > mp.mpf('0.40016')}")

# independent cross-check via mpmath's own digamma implementation
ref = mp.re(mp.digamma(mp.mpf('0.25')+7j))/2 - logpi/2
print(f"independent reference theta'(14) = {ref}")
print(f"reference inside certified interval? {theta_prime_lo <= ref <= theta_prime_hi}")
\end{lstlisting}

\textbf{Output:}
\begin{lstlisting}
S_199                 = 2.52604703934679486754882124953
tail bound            = 0.00376884422110552763819095477387
Re psi(1/4+7i) in     [1.94506253022415647930411820468, 1.95260021866636753458050011422]
theta'(14) certified  [0.400166322187378152580345426661, 0.403935166408483680218536381435]
exceeds 0.40016?      True
independent reference theta'(14) = 0.40048374392522506487181464471
reference inside certified interval? True
\end{lstlisting}

\subsection{Exact zero and transport identity for \texorpdfstring{$g(s)=s-s_0$}{g(s)=s-s0}}
\label{app:A2}

\begin{lstlisting}[language=Python]
"""
Appendix A.2 -- g(s)=s-s0 exact zero and transport-equation verification.
Reproduces eq:g-numerical-parameters, eq:g-numerical-zero,
and the h=1e-8 finite-difference check of the transport identity.
"""
import mpmath as mp
mp.mp.dps = 25

g1 = mp.im(mp.zetazero(1))
g2 = mp.im(mp.zetazero(2))
theta_g1 = mp.siegeltheta(g1)
theta_g2 = mp.siegeltheta(g2)
Omega1 = theta_g2 - theta_g1

def t_of_q(q):
    target = theta_g1 + Omega1*mp.mpf(q)
    t = g1 + q*(g2 - g1)
    for _ in range(8):
        f = mp.siegeltheta(t) - target
        fp = mp.mpf('0.5')*mp.log(t/(2*mp.pi))
        t = t - f/fp
    return t

def a_of_q(q):
    t = t_of_q(q)
    thetap = mp.mpf('0.5')*mp.log(t/(2*mp.pi))
    return thetap/Omega1

beta0 = mp.mpf('0.7')
q0 = mp.mpf('0.5')
s0 = beta0 + 1j*t_of_q(q0)
print(f"gamma_1={g1}, gamma_2={g2}, Omega_1={Omega1}")
print(f"s0 = {s0}")

def g(s):
    return s - s0

def G(beta, q):
    return g(beta + 1j*t_of_q(q))

print(f"g(0.7+i*t_1(1/2)) = {G(beta0, q0)}   (must be exactly 0)")

h = mp.mpf('1e-8')
dG_dbeta = (G(beta0+h, q0) - G(beta0-h, q0)) / (2*h)
dG_dq = (G(beta0, q0+h) - G(beta0, q0-h)) / (2*h)
lhs = -1j*a_of_q(q0)*dG_dq
print(f"dG/dbeta                = {complex(dG_dbeta)}")
print(f"-i*a(q)*dG/dq            = {complex(lhs)}")
print(f"residual (discretization only) = {complex(dG_dbeta - lhs)}")
\end{lstlisting}

\textbf{Output:}
\begin{lstlisting}
gamma_1=14.13472514173469379045725, gamma_2=21.02203963877155499262848, Omega_1=3.520041512807475232558556
s0 = (0.7 + 17.90557595350376093245695j)
g(0.7+i*t_1(1/2)) = (0.0 + 0.0j)   (must be exactly 0)
dG/dbeta                = (1+0j)
-i*a(q)*dG/dq            = (1.0001241821813218+0j)
residual (discretization only) = (-0.0001241821813218947+0j)
\end{lstlisting}
The residual is of order $h=10^{-8}$, confirming it is a finite-difference
discretization artifact rather than a failure of the exact identity
$\partial_\beta G=-ia(q)\partial_qG$, which holds symbolically for every
holomorphic $g$ by the transport identity of
Section~\ref{sec:universal-obstruction}.

\subsection{No-finite-jet perturbation}
\label{app:A3}

\begin{lstlisting}[language=Python]
"""
Appendix A.3 -- No-finite-jet perturbation confirmation (Proposition 14.1 / Remark 14.2).
Reproduces the 46.9 vs 49.0 comparison confirming the obstruction persists
after matching N=6 derivatives at the zero.
"""
import sympy as sp

qv = sp.Symbol('q')
qg = sp.Rational(1, 3)

# a concrete real-analytic F_beta with a simple zero at q=q_gamma
Fbeta = (qv - qg)*sp.sin(3*qv + 1) + (qv - qg)**2*sp.exp(qv)
# an unrelated reference F_{1/2,n}
Fhalf = sp.cos(2*qv) + qv**2

# best-fit C matching fourth derivatives at q_gamma
C_val = sp.diff(Fbeta, qv, 4).subs(qv, qg) / sp.diff(Fhalf, qv, 4).subs(qv, qg)
D4_original = sp.simplify(sp.diff(Fbeta - C_val*Fhalf, qv, 4))
probe_point = qg + sp.Rational(1, 10)
val_original = float(D4_original.subs(qv, probe_point))
print(f"D_q^4(F_beta - C*F_half) at probe point, BEFORE perturbation: {val_original}")

# perturb: agree with F_beta to order N=6 at q_gamma, arbitrary G elsewhere
G = sp.sin(5*qv) + 2
N = 6
Fbeta_tilde = Fbeta + (qv - qg)**(N+1)*G

print("\nDerivative agreement check (orders 0..N, must all be 0):")
for j in range(N+1):
    diff_j = sp.simplify(sp.diff(Fbeta_tilde - Fbeta, qv, j).subs(qv, qg))
    print(f"  order {j}: {diff_j}")

C_val_tilde = sp.diff(Fbeta_tilde, qv, 4).subs(qv, qg) / sp.diff(Fhalf, qv, 4).subs(qv, qg)
D4_tilde = sp.simplify(sp.diff(Fbeta_tilde - C_val_tilde*Fhalf, qv, 4))
val_tilde = float(D4_tilde.subs(qv, probe_point))
print(f"\nD_q^4(F_beta_tilde - C_tilde*F_half) at same point, AFTER perturbation: {val_tilde}")
print(f"Still generic/nonzero despite sharing 6 derivatives? {abs(val_tilde) > 1e-6}")
\end{lstlisting}

\textbf{Output:}
\begin{lstlisting}
D_q^4(F_beta - C*F_half) at probe point, BEFORE perturbation: 46.93821643404021

Derivative agreement check (orders 0..N, must all be 0):
  order 0: 0
  order 1: 0
  order 2: 0
  order 3: 0
  order 4: 0
  order 5: 0
  order 6: 0

D_q^4(F_beta_tilde - C_tilde*F_half) at same point, AFTER perturbation: 49.02760824764299
Still generic/nonzero despite sharing 6 derivatives? True
\end{lstlisting}
This example uses a generic real-analytic $F_\beta$, not
$\zeta(\beta+it_n(\cdot))$ itself, in accordance with the scope of
Proposition~\ref{prop:no-finite-jet} as clarified after
Corollary~\ref{cor:zero-trigger-not-finite-order}.

\begin{thebibliography}{99}

\bibitem{Ahlfors1979}
L. V. Ahlfors,
\emph{Complex Analysis}, 3rd ed.,
McGraw--Hill, New York, 1979.

\bibitem{AnseloneKorevaar1964}
P. M. Anselone and J. Korevaar,
\emph{Translation invariant subspaces of finite dimension},
Proc. Amer. Math. Soc. \textbf{15} (1964), 747--752.
DOI: 10.1090/S0002-9939-1964-0169048-7.

\bibitem{BaezDuarte2000}
L. B\'aez-Duarte, M. Balazard, B. Landreau and E. Saias,
\emph{Notes sur la fonction $\zeta$ de Riemann, 3},
Adv. Math. \textbf{149} (2000), no. 1, 130--144.
DOI: 10.1006/aima.1999.1861.

\bibitem{BaezDuarte2002}
L. B\'aez-Duarte,
\emph{New versions of the Nyman--Beurling criterion for the Riemann hypothesis},
Int. J. Math. Math. Sci. \textbf{31} (2002), no. 7, 387--406.
DOI: 10.1155/S0161171202013248.

\bibitem{BaezDuarte2003}
L. B\'aez-Duarte,
\emph{A strengthening of the Nyman--Beurling criterion for the Riemann hypothesis},
Rend. Lincei Mat. Appl. \textbf{14} (2003), no. 1, 5--11.

\bibitem{Burnol2002}
J.-F. Burnol,
\emph{A lower bound in an approximation problem involving the zeros of the Riemann zeta function},
Adv. Math. \textbf{170} (2002), no. 1, 56--70.
DOI: 10.1006/aima.2001.2066.

\bibitem{Beurling1955}
A. Beurling,
\emph{A closure problem related to the Riemann zeta-function},
Proc. Natl. Acad. Sci. USA \textbf{41} (1955), no. 5, 312--314.
DOI: 10.1073/pnas.41.5.312.

\bibitem{BombieriLagarias1999}
E. Bombieri and J. C. Lagarias,
\emph{Complements to Li's criterion for the Riemann hypothesis},
J. Number Theory \textbf{77} (1999), no. 2, 274--287.
DOI: 10.1006/jnth.1999.2392.

\bibitem{ChiangFeng2006}
Y.-M. Chiang and S.-J. Feng,
\emph{Difference independence of the Riemann zeta function},
Acta Arith. \textbf{125} (2006), no. 4, 317--329.
DOI: 10.4064/aa125-4-2.

\bibitem{Cicek2021}
F. \c{C}i\c{c}ek,
\emph{On the logarithm of the Riemann zeta-function near the nontrivial zeros},
Trans. Amer. Math. Soc. \textbf{374} (2021), 5995--6037.
DOI: 10.1090/tran/8426.

\bibitem{DLMF}
F. W. J. Olver, A. B. Olde Daalhuis, D. W. Lozier, B. I. Schneider,
R. F. Boisvert, C. W. Clark, B. R. Miller, B. V. Saunders,
H. S. Cohl and M. A. McClain, eds.,
\emph{NIST Digital Library of Mathematical Functions},
Chs. 5 and 25, National Institute of Standards and Technology.

\bibitem{Edwards1974}
H. M. Edwards,
\emph{Riemann's Zeta Function},
Academic Press, New York, 1974; Dover reprint, 2001.

\bibitem{Engert1970}
M. Engert,
\emph{Finite dimensional translation invariant subspaces},
Pacific J. Math. \textbf{32} (1970), no. 2, 333--343.
DOI: 10.2140/pjm.1970.32.333.

\bibitem{Ford2002}
K. Ford,
\emph{Zero-free regions for the Riemann zeta function},
in \emph{Number Theory for the Millennium II},
A K Peters, Natick, MA, 2002, 25--56.

\bibitem{GarunkstisLaurincikas2018}
R. Garunk\v{s}tis and A. Laurin\v{c}ikas,
\emph{Discrete mean square of the Riemann zeta-function over imaginary parts of its zeros},
Period. Math. Hungar. \textbf{76} (2018), no. 2, 217--228.
DOI: 10.1007/s10998-017-0228-6.

\bibitem{GarunkstisLaurincikasMacaitiene2017}
R. Garunk\v{s}tis, A. Laurin\v{c}ikas and R. Macaitien\.e,
\emph{Zeros of the Riemann zeta-function and its universality},
Acta Arith. \textbf{181} (2017), no. 2, 127--142.
DOI: 10.4064/aa8583-5-2017.

\bibitem{GarunkstisNovikas2025}
R. Garunk\v{s}tis and A. Novikas,
\emph{Discrete mean value theorems for the Riemann zeta-function over its shifted nontrivial zeros},
Funct. Approx. Comment. Math. \textbf{72} (2025), no. 2, 201--224.
DOI: 10.7169/facm/240929-7-10.

\bibitem{GarunkstisPaliulionyte2025}
R. Garunk\v{s}tis and J. Paliulionyt\.e,
\emph{Mean values of the Riemann zeta function at shifted zeros under the Riemann Hypothesis},
arXiv:2512.03297 (2025).

\bibitem{Gonek1984}
S. M. Gonek,
\emph{Mean values of the Riemann zeta-function and its derivatives},
Invent. Math. \textbf{75} (1984), no. 1, 123--141.
DOI: 10.1007/BF01403094.

\bibitem{Hardy1914}
G. H. Hardy,
\emph{Sur les z\'eros de la fonction $\zeta(s)$ de Riemann},
C. R. Acad. Sci. Paris \textbf{158} (1914), 1012--1014.

\bibitem{Ivic1985}
A. Ivi\'c,
\emph{The Riemann Zeta-Function: Theory and Applications},
Wiley, New York, 1985.

\bibitem{KaczorowskiPerelli1998}
J. Kaczorowski and A. Perelli,
\emph{Functional independence of the singularities of a class of Dirichlet series},
Amer. J. Math. \textbf{120} (1998), no. 2, 289--303.

\bibitem{KorolevLaurincikas2025}
M. Korolev and A. Laurin\v{c}ikas,
\emph{Joint functional independence of the Riemann zeta-function},
Indian J. Pure Appl. Math. \textbf{56} (2025), no. 4, 1523--1528.
DOI: 10.1007/s13226-024-00585-5.

\bibitem{Laurincikas2019}
A. Laurin\v{c}ikas,
\emph{Non-trivial zeros of the Riemann zeta-function and joint universality theorems},
J. Math. Anal. Appl. \textbf{475} (2019), no. 1, 385--402.
DOI: 10.1016/j.jmaa.2019.02.047.

\bibitem{LaurincikasSiauciunas2020}
A. Laurin\v{c}ikas and D. \v{S}iau\v{c}i\=unas,
\emph{Universality in short intervals of the Riemann zeta-function twisted by non-trivial zeros},
Mathematics \textbf{8} (2020), no. 11, 1936.
DOI: 10.3390/math8111936.

\bibitem{LevinsonMontgomery1974}
N. Levinson and H. L. Montgomery,
\emph{Zeros of the derivatives of the Riemann zeta-function},
Acta Math. \textbf{133} (1974), 49--65.
DOI: 10.1007/BF02392141.

\bibitem{Li1997}
X.-J. Li,
\emph{The positivity of a sequence of numbers and the Riemann hypothesis},
J. Number Theory \textbf{65} (1997), no. 2, 325--333.
DOI: 10.1006/jnth.1997.2137.

\bibitem{Lu2024}
F. L\"u,
\emph{On the algebraic difference independence of $\Gamma$ and $\zeta$},
J. Math. Anal. Appl. \textbf{538} (2024), no. 2, 128446.
DOI: 10.1016/j.jmaa.2024.128446.

\bibitem{MossinghoffTrudgianYang2024}
M. J. Mossinghoff, T. S. Trudgian and A. Yang,
\emph{Explicit zero-free regions for the Riemann zeta-function},
Res. Number Theory \textbf{10} (2024), article 11.
DOI: 10.1007/s40993-023-00498-y.

\bibitem{Nyman1950}
B. Nyman,
\emph{On Some Groups and Semigroups of Translations},
Ph.D. thesis, Uppsala University, 1950.

\bibitem{Ostrowski1920}
A. Ostrowski,
\emph{\"Uber Dirichletsche Reihen und algebraische Differentialgleichungen},
Math. Z. \textbf{8} (1920), 241--298.
DOI: 10.1007/BF01206530.

\bibitem{PlattTrudgian2021}
D. Platt and T. Trudgian,
\emph{The Riemann hypothesis is true up to $3\cdot10^{12}$},
Bull. Lond. Math. Soc. \textbf{53} (2021), no. 3, 792--797.
DOI: 10.1112/blms.12460.

\bibitem{Riemann1859}
B. Riemann,
\emph{\"Uber die Anzahl der Primzahlen unter einer gegebenen Gr\"osse},
Monatsber. K\"onigl. Preuss. Akad. Wiss. Berlin (1859), 671--680.

\bibitem{Speiser1934}
A. Speiser,
\emph{Geometrisches zur Riemannschen Zetafunktion},
Math. Ann. \textbf{110} (1934), 514--521.
DOI: \url{https://doi.org/10.1007/BF01448042}.

\bibitem{Szekelyhidi1982}
L. Sz\'ekelyhidi,
\emph{Note on exponential polynomials},
Pacific J. Math. \textbf{103} (1982), no. 2, 583--599.
DOI: 10.2140/pjm.1982.103.583.

\bibitem{Titchmarsh1986}
E. C. Titchmarsh, revised by D. R. Heath-Brown,
\emph{The Theory of the Riemann Zeta-Function}, 2nd ed.,
Clarendon Press, Oxford, 1986.

\bibitem{vonKoch1901}
H. von Koch,
\emph{Sur la distribution des nombres premiers},
Acta Math. \textbf{24} (1901), 159--182.
DOI: 10.1007/BF02403071.

\bibitem{Voronin1973}
S. M. Voronin,
\emph{On differential independence of $\zeta$-functions},
Dokl. Akad. Nauk SSSR \textbf{209} (1973), no. 6, 1264--1266.

\bibitem{Weil1952}
A. Weil,
\emph{Sur les ``formules explicites'' de la th\'eorie des nombres premiers},
Comm. S\'em. Math. Univ. Lund, Tome Suppl. (1952), 252--265.

\end{thebibliography}

\end{document}
